\documentclass[11pt,a4paper,oneside,openany]{book}
\usepackage{fontspec}
\usepackage{polyglossia}
\setdefaultlanguage{russian}
\setotherlanguage{english}
\setmainfont{Liberation Serif}
\setsansfont{Liberation Sans}
\setmonofont[Scale=0.88]{Liberation Mono}
\usepackage[a4paper,top=26mm,bottom=28mm,left=28mm,right=28mm,headsep=8mm]{geometry}
\usepackage{xcolor}
\definecolor{ink}{HTML}{1F2A37}
\definecolor{accent}{HTML}{2F5D8A}
\definecolor{accentgreen}{HTML}{3D7A4F}
\definecolor{accentamber}{HTML}{9A6B2F}
\definecolor{accentred}{HTML}{8F3B3B}
\definecolor{paleblue}{HTML}{EEF3F8}
\definecolor{palegreen}{HTML}{EDF5EF}
\definecolor{paleamber}{HTML}{F8F2EA}
\definecolor{palered}{HTML}{F8EDED}
\definecolor{palegray}{HTML}{F3F4F6}
\definecolor{midgray}{HTML}{8A8F98}
\usepackage{amsmath}
\usepackage{booktabs}
\usepackage{array}
\usepackage{longtable}
\usepackage{enumitem}
\usepackage{titlesec}
\usepackage{fancyhdr}
\usepackage{tikz}
\usetikzlibrary{arrows.meta,positioning,calc,shapes.misc}
\usepackage{pgfgantt}
\usepackage[most]{tcolorbox}
\usepackage{listings}
\usepackage{needspace}
\usepackage{microtype}
\usepackage[colorlinks=true,linkcolor=accent,urlcolor=accent,citecolor=accent]{hyperref}

\setlength{\parindent}{0pt}
\setlength{\parskip}{6pt plus 1pt}
\linespread{1.06}
\setlength{\emergencystretch}{1.5em}

% headings
\titleformat{\part}[display]{\sffamily\bfseries\centering}{\fontsize{14}{16}\selectfont\color{midgray}ЧАСТЬ \thepart}{12pt}{\fontsize{28}{32}\selectfont\color{ink}}
\titleformat{\chapter}[display]{\sffamily\bfseries}{\fontsize{44}{48}\selectfont\color{accent}\thechapter}{4pt}{\fontsize{22}{26}\selectfont\color{ink}}
\titlespacing*{\chapter}{0pt}{-6pt}{20pt}
\titleformat{\section}{\sffamily\bfseries\fontsize{13.5}{16}\selectfont\color{ink}}{\thesection}{0.6em}{}
\titlespacing*{\section}{0pt}{16pt plus 3pt minus 2pt}{6pt}
\titleformat{\subsection}{\sffamily\bfseries\fontsize{11.5}{14}\selectfont\color{ink}}{\thesubsection}{0.6em}{}
\titlespacing*{\subsection}{0pt}{12pt plus 2pt minus 2pt}{4pt}

% headers
\pagestyle{fancy}
\fancyhf{}
\renewcommand{\headrulewidth}{0.3pt}
\renewcommand{\chaptermark}[1]{\markboth{#1}{}}
\fancyhead[L]{\sffamily\footnotesize\color{midgray}\leftmark}
\fancyhead[R]{\sffamily\footnotesize\color{midgray}\thepage}
\fancypagestyle{plain}{\fancyhf{}\renewcommand{\headrulewidth}{0pt}\fancyfoot[C]{\sffamily\footnotesize\color{midgray}\thepage}}

% lists
\setlist[enumerate]{leftmargin=2.2em,itemsep=2pt,topsep=4pt,parsep=0pt}
\setlist[itemize]{leftmargin=1.6em,itemsep=2pt,topsep=4pt,parsep=0pt,label=--}

% boxes
\tcbset{
  basebox/.style={enhanced,breakable,boxrule=0pt,leftrule=3pt,arc=0pt,left=9pt,right=9pt,top=7pt,bottom=7pt,
    fonttitle=\sffamily\bfseries,attach title to upper={\par\vspace{2pt}},before skip=8pt,after skip=8pt}
}
\newtcolorbox{likbez}[1][Ликбез]{basebox,colback=paleblue,colframe=accent,coltitle=accent,title={#1}}
\newtcolorbox{whybox}[1][Почему так]{basebox,colback=paleamber,colframe=accentamber,coltitle=accentamber,title={#1}}
\newtcolorbox{invbox}[1][Инвариант]{basebox,colback=palegreen,colframe=accentgreen,coltitle=accentgreen,title={#1}}
\newtcolorbox{warnbox}[1][Осторожно]{basebox,colback=palered,colframe=accentred,coltitle=accentred,title={#1}}
\newtcolorbox{exbox}[1][Пример]{basebox,colback=palegray,colframe=midgray,coltitle=ink,title={#1}}

% listings
\lstset{basicstyle=\ttfamily\small,columns=fullflexible,keepspaces=true,breaklines=true,
  frame=none,backgroundcolor=\color{palegray},xleftmargin=8pt,xrightmargin=8pt,framexleftmargin=8pt,
  aboveskip=8pt,belowskip=8pt,showstringspaces=false}

% tikz
\tikzset{
  box/.style={draw=ink,rectangle,thick,align=center,inner sep=6pt,font=\small,fill=white},
  boxa/.style={box,fill=paleblue,draw=accent},
  boxg/.style={box,fill=palegreen,draw=accentgreen},
  boxw/.style={box,fill=paleamber,draw=accentamber},
  lbl/.style={font=\footnotesize,align=center,color=ink},
  arr/.style={-{Latex[length=2.2mm]},thick,color=ink},
  darr/.style={{Latex[length=2.2mm]}-{Latex[length=2.2mm]},thick,color=ink},
}

\newcommand{\term}[1]{\textbf{#1}}
\newcommand{\inv}[1]{\textbf{I#1}}

\begin{document}
\frontmatter
\begin{titlepage}
\thispagestyle{empty}
\vspace*{30mm}
{\sffamily\bfseries\fontsize{34}{40}\selectfont\color{ink}Рой над единой средой\par}
\vspace{10pt}
{\sffamily\fontsize{16}{20}\selectfont\color{midgray}Руководство инженера по протоколу\\распределённой агентной разработки\par}
\vspace{22mm}
{\normalsize\color{ink}
Что это за проект, какие проблемы он решает,\\
из каких примитивов собирается\\
и с чего начинается его реализация\par}
\vfill
{\normalsize\color{ink}
По материалам работы\\
\textit{Разработка программного обеспечения роем автономных агентов над единой средой исполнения}\\
Алексей Григорьев, Иван Опарин, Relux Works, третья редакция, версия 3.3\par}
\vspace{8mm}
{\sffamily\footnotesize\color{midgray}Черновик 0.4, сентябрь 2026 \quad github.com/relux-works\par}
\end{titlepage}

\chapter*{Как читать эту книгу}
\addcontentsline{toc}{chapter}{Как читать эту книгу}

Эта книга написана для инженера, который присоединяется к проекту. Она предполагает десять и больше лет опыта в программировании, хорошее понимание того, как устроены современные агентские системы поверх больших языковых моделей, и не предполагает опыта с распределёнными базами данных. Всё, что нужно знать о репликации, кворумах, арендах и деревьях Меркла, объясняется по ходу в отдельных врезках.

Книга состоит из трёх частей. Первая отвечает на вопрос зачем: какую проблему мы решаем и почему существующие инструменты её не решают. Вторая описывает, как устроен протокол: среда исполнения, журнал, полномочия, оркестратор, цикл задачи, рабочие ноды, модель угроз. Третья описывает, как это построить: из каких существующих компонентов собирается система, что нужно написать заново, как выглядит план на двенадцать недель и какие вопросы остаются открытыми.

Источник для книги это третья редакция работы (версия 3.3), на которой основан протокол, два независимых ревью, которые её сформировали, и обсуждения, которые к ней привели. Там, где книга уходит в детали реализации глубже, чем работа, это отмечено. Там, где решение ещё не принято, это сказано прямо.

Версия 0.4 выравнивает книгу с системой Apiary на 29 сентября 2026 года. Протокол книги есть распределённый слой этой системы, и большинство его понятий теперь принадлежат спецификациям модулей; где книга и спецификация расходятся, права спецификация. Карта разделов работы на их нынешних владельцев дана во введении работы. Главы 10--12 заменены роадмапом экосистемы и репозиториями модулей и сохранены как история, открытые вопросы главы 13 указывают, где вопрос живёт сейчас, а карта имён для людей вынесена в глоссарий.

В тексте используются пять видов врезок.

\begin{likbez}
Объяснение понятия из распределённых систем или криптографии, которое нужно для понимания следующего раздела. Если вы знакомы с темой, врезку можно пропустить.
\end{likbez}

\begin{whybox}
Обоснование архитектурного решения: какие альтернативы были и почему выбрана именно эта.
\end{whybox}

\begin{invbox}
Формулировка одного из десяти инвариантов протокола. Спецификации нижнего уровня ссылаются на инварианты по номерам I1 и далее.
\end{invbox}

\begin{warnbox}
Место, где легко ошибиться при реализации, или свойство, которое протокол не обещает.
\end{warnbox}

\begin{exbox}
Конкретный пример: сценарий, последовательность записей, набросок структуры данных.
\end{exbox}

Терминология выровнена с работой и с существующими спецификациями компонентов. Одно слово, эпоха, в разных компонентах означает разные вещи; в этой книге оно всегда означает срок аренды оркестратора, а остальные значения оговорены там, где встречаются. Глоссарий вынесен в приложение.

\tableofcontents
\mainmatter
\part{Зачем}
\chapter{Проблема}

\section{Как сегодня устроена агентная разработка}

Почти все инструменты, в которых агенты пишут код, собраны вокруг одной схемы. Агент запускается там, где лежит код. Изоляция параллельной работы достигается копиями рабочего дерева: git worktree на локальной машине или отдельная виртуальная машина в облаке. Параллелизм достигается умножением инстансов: десять задач, десять агентов, десять копий среды.

Для веб-сервисов и бэкендов эта схема работает хорошо. Среда там дешёвая: контейнер с компилятором и зависимостями поднимается за секунды, и облачные песочницы на каждого агента стоят копейки. Ровно поэтому индустрия оптимизирует именно эту нишу, и ровно поэтому она не смотрит на другие.

Схема наследует свою главную предпосылку у git. Git создавался для людей и предполагает, что у каждого разработчика есть собственная копия среды. Для человека это разумно: людей мало, их машины сильные, и каждый работает в своём темпе. Для агентов та же предпосылка превращается в три проблемы.

\section{Три проблемы}

\subsection{Среда дорогая или невоспроизводимая}

Сборка iOS-приложения требует macOS и Xcode конкретной версии. Подпись требует ключей, которые нельзя раздавать. Тесты на устройствах требуют самих устройств. Обучение или инференс модели требует GPU, которых нет в дешёвой песочнице. Работа с данными компании требует, чтобы данные не покидали её периметр. Большой монорепозиторий требует часов на прогрев кешей сборки.

Во всех этих случаях воспроизвести среду на каждого агента либо очень дорого, либо невозможно. Схема с копией среды на инстанс здесь просто не работает, и команды остаются с одним агентом на одной машине.

\subsection{Лимиты провайдера на одного владельца}

Подписка на модель выдаёт квоту одному владельцу: столько-то токенов в окно времени. Десять агентов под одной подпиской делят одну квоту, и рой упирается в лимит раньше, чем в вычислительную мощность. Обойти это можно только большим числом подписок, и здесь начинается третья проблема.

\subsection{Чужие агенты на одной машине означают чужие учётные данные на ней}

Пусть над проектом работают пять человек, у каждого своя подписка. Чтобы их агенты работали на одной машине с дорогой средой, каждый должен передать этой машине свои учётные данные. Условия провайдеров это запрещают, а здравый смысл не советует: машина, на которой лежат чужие ключи, становится целью. Значит, объединить подписки нескольких людей на одной среде сегодня нельзя, не нарушив либо правила, либо безопасность.

\section{Неравенство}

Из трёх проблем следует неравенство, которое и есть настоящий предмет проекта.

Большие рои агентов, координируемые для работы над одним проектом, доступны сегодня вендорам моделей, владельцам крупных вычислительных ресурсов и корпорациям. У них есть и квота, и инфраструктура, и возможность поставить среду рядом с агентами. Открытые проекты активно используют агентов, но спроектированы под координацию через людей: issue, pull request, ревью мейнтейнера. Агент в открытом проекте это помощник одного человека, а не участник роя.

Мы предполагаем, что координация на уровне автономного роя эффективнее координации через людей: человек читает pull request часами, рой обменивается назначениями и вердиктами за секунды. Это гипотеза, и глава 11 описывает эксперимент, который её проверяет; в работе она нигде не утверждается как факт. Разрыв между теми, кто может собрать рой, и всеми остальными в любом случае растёт вместе с возможностями моделей.

Проект есть ответ на это неравенство: протокол, который позволяет незнакомым друг с другом людям приносить в общий проект не патчи, а ёмкость своих подписок и моделей, и делает это безопасным для проекта и для них самих.

\section{Что мы утверждаем и чего не утверждаем}

Важно с самого начала провести границу честно, потому что от неё зависит, как оценивать результат.

Мы не утверждаем, что протокол делает агентов умнее. Качество декомпозиции задач, вероятность того, что ревьювер заметит дефект, и сила тестового набора остаются свойствами моделей и проекта. Протокол ни одну из этих трёх величин не улучшает.

Мы утверждаем другое, и формулировка вклада в третьей редакции работы звучит так: мы рассматриваем совместное изменение проекта агентами независимых владельцев в домене, где у принятия результата нет механического оракула; рабочая нода может предлагать и производить изменения, но не может расширить область задачи, изменить условия её приёмки или продвинуть основную ветку; авторизация, исполнение и принятие связаны проверяемыми записями, относящимися к конкретным версиям контракта задачи и состояния проекта. Иначе говоря, никакой агент не получает больше, чем ему выдано, и ничего не попадает в основную ветку без независимой проверки ровно того состояния, которое попадает. Протокол делает суждение агентов независимым, случайным и подотчётным: ревьювер выбран так, что его нельзя подобрать, и каждое действие оставляет след, по которому его можно воспроизвести. Протокол делает возможным то, что сегодня доступно немногим: рой из ёмкости независимых людей.

Ценность здесь это доступ и контроль. Эффективность нужно измерять, а не декларировать: принятые задачи на миллион токенов, время от постановки до интеграции, доля откатов, часы человеческого ревью на принятую задачу. Эти метрики встроены в систему с первого дня, и глава 12 объясняет, как они собираются.

\section{Сценарий}

Чтобы дальнейший текст не был абстрактным, зафиксируем сценарий, к которому книга будет возвращаться.

\begin{exbox}[Сценарий: приложение с нативным ядром]
Проект: iOS-приложение с ядром на Rust. Среда исполнения: один Mac Studio с Xcode, ключами подписи и подключёнными устройствами, стоящий у кор-команды. Кор-команда: два человека, у каждого своя кор-нода. Контрибьюторы: пять человек в разных городах, у троих подписки на разные модели, у одного ключ API, у одного локальная модель на собственном Mac. Никто из них не имеет доступа к Mac Studio, ключам подписи или устройствам.

Кор-команда ставит цель: перенести модуль синхронизации на новый протокол. Оркестратор, одна из кор-нод, раскладывает цель на двадцать задач с областями изменений и критериями приёмки. Каждую задачу подхватывает агент одного из контрибьюторов, работает в изолированном каталоге на Mac Studio через удалённый shell, представляет коммит. Случайно выбранный агент другого контрибьютора делает ревью. Тесты гоняет третий. Среда сливает результат в основную ветку только тогда, когда журнал подтвердил, что ревью и тесты относятся ровно к тому, что сливается. Если что-то идёт не так, любая кор-нода останавливает всё одной подписью.
\end{exbox}

Всё, что дальше, есть объяснение того, как этот сценарий работает и почему его нельзя сломать, не сломав явных допущений.

\section{Контекст: что уже сделали другие}

Объединение чужих агентов на общем проекте само по себе уже не ново, и это важно понимать до чтения остального. В формальной математике открытая платформа Prove2Me позволяет любому владельцу агента вносить доказательства в общую миссию: граф формулировок задаёт фронт работы, формулировки отделены от доказательств, описание каждого узла делает результаты находимыми и переиспользуемыми. На ней в сентябре 2026 года была завершена формализация Великой теоремы Ферма силами десятков агентов, а в отдельном эксперименте три личные подписки на модель за три дня совместно формализовали крупный результат. Для обычной разработки ближайшая работа это команда из шестнадцати параллельных агентов, построившая компилятор Си: агенты в отдельных контейнерах, общий репозиторий, блокировки задач, самостоятельная интеграция.

Чем мы отличаемся. В формальной математике у принятия есть оракул, ядро проверки доказательств, и вся координация стоит на нём. У кода оракула нет, и принятие приходится конструировать: отсюда независимое ревью, кандидат, аттестации и кворум. В эксперименте с компилятором все агенты одного владельца и доверены; у нас нет. Наш контракт задачи и граф зависимостей есть перенос разделения формулировок и доказательств на код, а всё остальное есть ответ на отсутствие оракула и на недоверие к участникам.

\chapter{Идея целиком}

\section{Три плоскости}

Проще всего увидеть архитектуру как три плоскости, которые в обычной агентной системе склеены в одном процессе, а здесь разведены по разным машинам и разным владельцам.

\begin{center}
\begin{tikzpicture}[node distance=6mm]
\node[boxa,minimum width=134mm,minimum height=13mm] (think) {Плоскость мышления\\[1pt]\footnotesize агентские циклы на машинах контрибьюторов: модель, контекст, ключи, ёмкость};
\node[boxg,minimum width=134mm,minimum height=13mm,below=of think] (ctrl) {Плоскость управления\\[1pt]\footnotesize кор-ноды: журнал записей, правила, полномочия, оркестратор};
\node[boxw,minimum width=134mm,minimum height=13mm,below=of ctrl] (exec) {Плоскость исполнения\\[1pt]\footnotesize среда проекта: рабочее дерево, санбоксы, тесты, основная ветка};
\draw[darr] (think) -- node[lbl,right] {назначения, ревизии, вердикты} (ctrl);
\draw[arr] (ctrl) -- node[lbl,right] {записи с сертификатом кворума} (exec);
\draw[arr] ($(think.west)+(-4mm,0)$) -- ++(-6mm,0) |- node[lbl,left,pos=0.25,rotate=90,anchor=south] {удалённый shell в санбоксе задачи} ($(exec.west)+(-4mm,0)$);
\end{tikzpicture}
\end{center}

Плоскость мышления это агентские циклы: модель, её контекст, инструменты, которые она вызывает. Они работают на машинах контрибьюторов под их собственными учётными данными и с их собственной ёмкостью. Плоскость исполнения это среда проекта: единственная машина или небольшой пул машин, где лежит код, собираются артефакты и гоняются тесты. Плоскость управления это журнал записей и кор-ноды, которые его ведут: кто что вправе делать, что назначено, что проверено, что интегрировано.

Между плоскостями проходят только команды, их вывод, хеши и подписи. Учётные данные не покидают машину контрибьютора. Код не покидает среду проекта иначе как в санбоксе конкретной задачи. Полномочия существуют только в журнале.

\section{Инверсия: среда одна, агенты снаружи}

Обычная схема: агент запускается там, где код. Наша: код лежит в одном месте, а агенты подключаются к нему снаружи. Среда не запускает агентских циклов вообще. Она пассивна: исполняет команды, поступающие извне, и только в границах, заданных журналом.

Это инверсия направления, и у неё есть цена и выигрыш. Цена: среда становится единственной точкой, через которую проходит вся работа, и её ёмкость ограничивает размер роя. Выигрыш: среду не нужно воспроизводить, ключи не нужно раздавать, а квоты приносят сами участники. Для дорогих сред выигрыш перевешивает.

\begin{whybox}[Почему не облачная песочница на агента]
Облачные песочницы решают изоляцию, воспроизводя среду на каждого агента. Для Xcode, ключей подписи, устройств, GPU и данных с ограничением на вывоз это либо невозможно, либо дорого. Кроме того, песочница всё равно принадлежит одному аккаунту с одной квотой; объединение квот разных людей она не решает.
\end{whybox}

\section{Журнал полномочий вместо git как шины}

Git остаётся в системе на двух ролях: хранилище, адресуемое по содержимому (любой коммит имеет хеш, и по хешу можно проверить, что именно проверял ревьювер), и механизм ветвления внутри одной машины. Git перестаёт быть шиной координации между участниками. Координация идёт через журнал.

Журнал это упорядоченная цепочка подписанных записей: цель поставлена, задача создана, задача назначена такой-то ноде, ревизия с таким-то хешем представлена, ревью такое-то, тесты такие-то, интеграция такая-то. Каждая запись ссылается на предыдущую по хешу и подписана достаточным числом кор-нод. Среда исполняет только то, что подтверждено журналом: открывает санбокс по записи назначения, сливает в основную ветку по записи об интеграции.

Так всякое действие в системе имеет ответ на вопрос кто и на каком основании его разрешил, и ответ этот проверяем без доверия к тому, кто его даёт.

\section{Сессия это кэш, журнал это истина}

У агента есть контекст: история диалога, промежуточные выводы, план. В обычных системах этот контекст и есть состояние проекта: если процесс упал, состояние потеряно. Здесь состояние проекта живёт в журнале, а контекст модели есть кэш. Оркестратор, который потерял сессию, восстанавливается из журнала. Оркестратор, который получил сессию от предшественника, обязан считать её недоверенной и разрешать любое расхождение в пользу журнала.

Это одно из самых важных решений в проекте, и глава 7 объясняет, почему без него не работает ни смена оркестратора, ни защита от отравленного контекста.

\section{Что мы не верифицируем}

Через консенсус кор-нод проходят только детерминированные предикаты: допустим ли переход состояния, лежит ли область задачи внутри области цели, относятся ли ревью и тесты к тому же хешу, что интегрируется, выбран ли ревьювер по случайности, которую никто не контролировал. Кор-ноды не голосуют о том, хорош ли код и прав ли ревьювер.

Это принципиально. Криптография и консенсус умеют доказывать, что процедура соблюдена: кто, когда и на каком основании. Они не умеют доказывать, что модель правильно поняла задачу. Попытка проголосовать о правильности суждения ломает и консенсус, и суждение. Поэтому суждение остаётся снаружи, а протокол делает его независимым (продьюсер, ревьювер и тестировщик разные), случайным (ревьювера нельзя подобрать) и подотчётным (каждое действие оставляет след).

\section{Уровни доверия}

Протокол определён для трёх уровней доверия с одинаковыми инвариантами и разным определением того, сколько подписей нужно, чтобы запись считалась совершённой.

\begin{center}
\small
\begin{tabular}{p{1.5cm} p{3.3cm} p{3.1cm} p{5.2cm}}
\toprule
Уровень & Сколько подписей & Порядок записей & От чего защищает сертификат \\
\midrule
L0 & одна, единственная кор-нода & единственная нода & от действий вне кора; кор и среда доверены \\
L1 & больше половины из $n$ & доверенный секвенсер & от одиночного несанкционированного решения при компрометации меньшинства ключей; от раздвоения истории защищает только честность секвенсера \\
L2 & $2f+1$ из $n \geq 3f+1$ & византийски устойчивый консенсус & и от несанкционированного решения, и от раздвоения истории при $f$ злонамеренных кор-нодах \\
\bottomrule
\end{tabular}
\end{center}

Первое развёртывание идёт на L0: одна кор-нода, один подписант. Все механизмы при этом уже работают, просто кворум равен одной подписи. Повышение уровня не требует переписывать ничего, кроме определения кворума. Последний столбец важнее, чем кажется: читатель должен понимать, защищает ли сертификат от одиночного несанкционированного решения, от раздвоения истории или от обоих. Глава 3 объясняет, откуда берутся $2f+1$ и $3f+1$ и почему порог большинства на L1 не даёт единственности истории.

\section{Десять инвариантов}

Всё остальное в книге есть способы обеспечить десять инвариантов. Здесь они перечислены кратко, чтобы читатель видел каркас; глава 5 разбирает каждый и показывает, что ломается без него.

\begin{enumerate}[label=I\arabic*.,leftmargin=2.6em]
\item Ни одного эффекта на среде без совершённой и актуальной записи авторизации с совпадающим digest плана.
\item Полномочия только сужаются: от цели к задаче и от задачи к ревизии; условия приёмки меняются только новой версией контракта с авторизацией.
\item Основная ветка продвигается только атомарной заменой на неизменяемого кандидата, к которому адресованы финальное одобрение и аттестации.
\item Суждение не входит в консенсус.
\item Продьюсер, ревьювер и исполнитель проверок независимы согласно объявленной политике, которая не ослабляется автоматически; случайность выбора фиксируется после того, как подбирать уже нечего.
\item Ключи контрибьютора не покидают его машину; секреты проекта существуют только в раннере.
\item Журнал есть истина, сессия есть кэш; записи устаревшей аренды отвергаются.
\item Любая кор-нода останавливает, кворум возобновляет; откат журнала невозможен.
\item Ограничение есть сообщаемый факт; требующий полного ограничения отказывает на частичном.
\item Все обращения к среде через протокол имеют проверяемое происхождение; полнота внутреннего контекста недоверенной ноды не удостоверяется.
\end{enumerate}

Если вы запомните из книги только этот список и три плоскости, вы уже сможете отличить правильное решение от неправильного почти в любом вопросе реализации.

\part{Как это устроено}
\chapter{Примитивы, на которых всё стоит}

Эта глава для читателя, который писал большие системы, но не строил распределённых баз данных. Каждый раздел объясняет один примитив и говорит, зачем он нужен именно нам. Если примитив вам знаком, раздел можно пропустить; в остальных главах на них есть ссылки.

\section{Реплицированный журнал и порядок}

Самая простая распределённая структура данных это журнал: последовательность записей, к которой можно только дописывать. Если несколько машин держат одинаковую копию журнала, они держат одинаковое состояние, потому что состояние есть результат применения записей по порядку. Все сложные вещи в распределённых базах сводятся к одному вопросу: как договориться о порядке записей.

Договориться можно двумя способами. Первый: назначить одну машину секвенсером, который присваивает записям номера, а остальные лишь проверяют и подписывают. Второй: провести консенсус, где порядок фиксируется голосованием без единого назначающего. Первый способ проще и живее, но секвенсер может остановиться или соврать. Второй сложнее, но не имеет единой точки.

\begin{likbez}[Зачем нам]
Журнал протокола есть ровно такой реплицированный журнал: записи о целях, задачах, назначениях, ревью и интеграциях. На уровнях L0 и L1 порядок задаёт секвенсер, а кор-ноды подписывают; на L2 порядок фиксируется консенсусом. Разделение порядка и подписи нужно, чтобы журнал продвигался, даже когда кор-ноды это ноутбуки, которые бывают выключены: секвенсер это всегда включённый сервис проекта, а подписи приходят, когда ноутбуки в сети.
\end{likbez}

\section{Отказы: сбой и византийское поведение}

Распределённые протоколы различают два класса отказов. Сбой: машина остановилась или отвалилась от сети, но не врёт. Византийское поведение: машина делает что угодно, включая ложь, противоречивые сообщения разным собеседникам и сговор с другими. Название идёт из задачи о византийских генералах, которые должны договориться об атаке, когда часть из них предатели.

Протоколы для сбоев (Raft, Paxos) требуют, чтобы большинство машин было живо. Протоколы для византийского поведения (PBFT, Tendermint, HotStuff) требуют, чтобы злонамеренных было меньше трети.

\begin{likbez}[Откуда берётся треть]
Пусть машин $n$, злонамеренных не больше $f$. Решение принимается, когда собрано $q$ подписей. Мы не можем ждать всех, потому что $f$ машин могут просто молчать; значит, $q \leq n - f$. Среди собранных $q$ подписей до $f$ могут быть от лжецов, которые в другом месте подписали противоположное. Чтобы два противоречивых решения не могли оба набрать $q$, их честные подписанты должны пересекаться: $2(q - f) > n - f$, то есть $q > (n+f)/2$. Соединяя с $q \leq n - f$, получаем $n > 3f$, откуда $n \geq 3f + 1$ и $q = 2f + 1$. При четырёх машинах терпим одну лживую, при семи две.
\end{likbez}

\begin{likbez}[Зачем нам]
Кор-ноды на L2 могут быть скомпрометированы: взломанная машина, отравленный инжектом агент, недобросовестный участник. Требование $n \geq 3f+1$ и кворум $2f+1$ гарантируют, что ни одна запись не будет совершена дважды с противоречивым содержимым и что $f$ лжецов не подделают историю. На L0 и L1 византийская устойчивость не заявляется, и кворум проще.
\end{likbez}

\section{Кворум и сертификат кворума}

Кворум это достаточное число подписей под одним и тем же содержимым. Набор подписей, собранный под записью, называют сертификатом кворума, QC. Полезное свойство: QC можно передавать и проверять офлайн. Тот, кто получил запись с QC, не должен ни с кем связываться, чтобы убедиться, что запись совершена: достаточно проверить подписи по известным ключам.

\begin{likbez}[Зачем нам]
Среда исполнения проверяет QC на записи назначения, прежде чем открыть санбокс, и на записи об интеграции, прежде чем продвинуть основную ветку. Ей не нужно доверять ни оркестратору, ни секвенсеру: подписи говорят сами за себя. Определение кворума и есть параметр, который меняется между уровнями L0, L1 и L2.
\end{likbez}

\begin{likbez}[Чего порог большинства не даёт]
Пусть подписантов трое, порог два, и один из них, $A$, злонамерен. $A$ подписывает две несовместимые записи на одно и то же место в журнале; честный $B$ подписывает первую, честный $C$ вторую. Обе записи получают сертификат большинства: $\{A, B\}$ и $\{A, C\}$. Порог защитил от одиночного решения, но не от раздвоения истории. Пересечение двух кворумов размера $t$ содержит не меньше $2t - n$ участников, и чтобы в нём гарантированно был честный при $f$ злонамеренных, нужно $2t - n > f$; для $n = 3$, $t = 2$, $f = 1$ это не выполнено. Поэтому L1 у нас не обещает единственности истории силой подписей: порядок там задаёт доверенный секвенсер, а подписи контролируют авторизацию. Единственность истории при лжецах даёт только L2 с полноценным византийским консенсусом.
\end{likbez}

\section{Лидер, аренда и fencing-токен}

Во многих системах есть роль, которую в каждый момент должна держать ровно одна машина: лидер, владелец, оркестратор. Оформляют её арендой: записью о том, что машина X держит роль до момента T, и продлевают heartbeat. Если heartbeat пропал, аренда истекает и роль переходит другому.

Аренда сама по себе не защищает от классической ошибки. Машина X держит аренду, зависает на минуту (сборка мусора, сеть, сон ноутбука), аренда истекает, роль переходит Y, а потом X просыпается и, ничего не зная, продолжает работать как лидер. Теперь лидеров два.

Лекарство называется fencing-токен: с каждой арендой выдаётся номер, который только растёт, и каждое действие лидера несёт этот номер. Тот, кто исполняет действия, помнит наибольший виденный номер и отвергает всё, что пришло с меньшим. Проснувшийся X действует с устаревшим номером, и его действия отбрасываются.

\begin{likbez}[Зачем нам]
Оркестратор есть арендуемая роль. Каждая его запись несёт эпоху и fencing-токен аренды; валидаторы отвергают записи с устаревшим токеном, но сохраняют их в ветке расхождений для аудита. Так двое оркестраторов не могут одновременно раздавать задачи, даже если один из них об этом не знает. Тот же приём защищает передачу сессии: принудительный захват отсекает предшественника токеном.
\end{likbez}

\section{Хеш-цепочка, дерево Меркла и доказательство включения}

Хеш-цепочка: каждая запись содержит хеш предыдущей. Изменить запись в середине нельзя, не изменив все последующие хеши, а последний хеш, голову, все знают. Дерево Меркла: над хешами записей строится двоичное дерево, где каждый узел есть хеш двух детей, и корень дерева фиксирует весь журнал одним числом. Полезное свойство дерева: чтобы доказать, что запись входит в журнал с данным корнем, достаточно показать логарифм от длины журнала хешей, а не весь журнал. Это доказательство включения.

Снапшот это подписанный набор из корня, длины журнала, головы и версии. Кто знает снапшот, тот может проверить любую запись по доказательству включения, не имея журнала целиком.

\begin{likbez}[Зачем нам]
Среда исполнения есть светлый клиент журнала: ей не нужно хранить журнал целиком, достаточно последнего снапшота с QC и доказательства включения для конкретной записи назначения или интеграции. Журналы прозрачности, на которых стоит инфраструктура сертификатов и подписи артефактов, устроены ровно так.
\end{likbez}

\section{Защита от отката}

Снапшоты подписаны, но подпись не мешает подсунуть светлому клиенту старый, тоже подписанный снапшот. Если клиент примет старый снапшот, он согласится с состоянием, в котором ещё не было, скажем, отзыва ключа или остановки. Защита: клиент хранит наивысшую принятую версию журнала как долговременное состояние безопасности, а не как кэш, и отвергает всё, что старше. Это состояние переживает перезапуски и не сбрасывается при ротации ключей.

\begin{likbez}[Зачем нам]
Без защиты от отката сертификат экстренной остановки можно было бы обойти, подсунув среде снапшот до остановки. Инвариант I8 требует высшей принятой версии именно поэтому.
\end{likbez}

\section{Атомарная замена и очередь интеграции}

Compare-and-swap: замени значение на новое, но только если текущее равно ожидаемому. Если кто-то успел изменить значение раньше, операция не выполняется, и надо повторить с новым ожиданием. Это единственный способ безопасно продвигать одну общую ссылку, за которую конкурируют несколько писателей.

В системах интеграции кода из этого выросло правило: никогда не сливать состояние, которое не было собрано и протестировано ровно в том виде, в каком оно попадёт в основную ветку. Очередь интеграции берёт кандидатов по одному, строит слияние с текущей головой, тестирует именно его и только потом продвигает голову.

\begin{likbez}[Зачем нам]
Одобрение и проверки привязаны к хешу. Если ревьюили коммит $R$, а в основную ветку попадает слияние $R$ с ушедшей вперёд головой $M$, то попадает другой хеш и, возможно, другое поведение, даже если сам патч не изменился. Поэтому среда строит интеграционного кандидата $I$, финальное одобрение и проверки относятся к $I$, а голова продвигается атомарно с $M$ на $I$; если голова уже не $M$, кандидат перестраивается, и одобрение с проверками повторяются. Это инвариант I3.
\end{likbez}

\section{Доступность данных}

Журнал хранит хеши тяжёлых объектов, а не сами объекты: тела задач, находки ревью, снимки сессий. Хеш гарантирует целостность объекта, но не его наличие. Кор-ноды могут совершить запись с хешем объекта, который потом никто не сможет получить, потому что единственная копия лежала на выключенном ноутбуке.

\begin{likbez}[Зачем нам]
Перед включением хеша в журнал требуется сертификат доступности: несколько независимых реплик подтвердили, что хранят объект. Для чекпоинтов сессии оркестратора это критично: без чекпоинта преемник стартует холодно.
\end{likbez}

\section{Случайность, которую никто не контролирует}

Если ревьювера выбирать случайно, кто-то должен сгенерировать случайность. Если её генерирует оркестратор, он может подобрать удобную. Если она известна заранее, оркестратор может подобрать момент или состав кандидатов, чтобы случайность привела к нужному ревьюверу; это называется grinding.

Решения: случайность собирается из вкладов нескольких сторон так, что ни одна не контролирует итог (пороговый маяк, commit-reveal), либо вычисляется проверяемой случайной функцией из секрета одной стороны так, что результат нельзя подобрать, а проверить может любой. И в любом случае порядок такой: сначала фиксируется множество кандидатов и объект, для которого выбирают, и только потом получается случайность.

\begin{likbez}[Зачем нам]
Множество кандидатов в ревьюверы замораживается записью до появления хеша ревизии, случайность получается после, и ревьювер вычисляется из идентификатора задачи, хеша ревизии и случайности. Любой валидатор пересчитывает выбор. Это половина инварианта I5.
\end{likbez}

\section{Сибил-атака}

В открытой системе один участник может завести много личностей и получить больше голосов, больше назначений, больше репутации. Единственная универсальная защита это стоимость личности. Доказано, что не существует нетривиальной репутационной системы, которая одновременно устойчива к множественным личностям и вознаграждает пропорционально вкладу; можно только сделать атаку дорогой.

\begin{likbez}[Зачем нам]
Нода это ёмкость, и ёмкость стоит денег: подписка, ключ API, железо для локальной модели. Это естественный барьер, но не абсолютный, поэтому он дополняется репутацией и адаптивной глубиной проверки. Публичный рой не обещает пропорциональной справедливости, он обещает дорогую атаку.
\end{likbez}

\section{Что мы берём у готовых протоколов}

Raft даёт нам модель снимка и лога для передачи состояния оркестратора: преемник загружает последний чекпоинт и докатывает записи журнала после него. PBFT, Tendermint и HotStuff дают модель византийского консенсуса для L2, включая смену лидера консенсуса при его молчании или лжи. Журналы прозрачности дают форму снапшотов и доказательств включения. Очереди интеграции дают правило никогда не сливать непроверенное. Ничего из этого мы не изобретаем; наш вклад в композиции и в том, что предметом журнала являются полномочия агентов.

\chapter{Среда исполнения}

Среда исполнения это плоскость, где живёт код и где происходят все эффекты. Она единственная в логическом смысле и пассивная в поведенческом: не запускает агентских циклов, не принимает решений, исполняет только то, что подтверждено журналом. Эта глава описывает, из чего она состоит и почему каждая часть устроена именно так.

\section{Авторитетная среда и реплики}

Логически состояние проекта одно: основная ветка и журнал определяют его полностью. Физически среда может быть пулом машин. Одна из них авторитетна: только на ней продвигается голова основной ветки. Остальные, исполнительные реплики, разделяют рабочее дерево по снимкам и обслуживают санбоксы задач. Отдельная роль, верификатор, повторяет критические прогоны тестов независимо от той машины, где работал продьюсер.

\begin{whybox}
Единственная физическая машина была бы и узким местом, и единственной точкой отказа. Пул реплик снимает ограничение ёмкости, а верификатор даёт хотя бы частичную независимость тестов от машины, которую мы вынуждены считать доверенной. Тезис о дорогой среде при этом не ломается: дорого воспроизвести среду на каждого контрибьютора, а не на две-три машины кор-команды.
\end{whybox}

\section{Санбокс: копия рабочего дерева не является границей}

Санбокс задачи состоит из двух вещей. Первая: копия рабочего дерева на момент назначения, созданная копированием при записи. На macOS это клон файлов APFS, который создаётся за миллисекунды и не занимает места, пока файл не изменён; на Linux снимок btrfs или ZFS либо overlay. Вторая: политика ограничения на уровне операционной системы.

\begin{warnbox}[Worktree это не безопасность]
Первая редакция протокола описывала санбокс как отдельный каталог с ограничением путей. Этого недостаточно. Копия рабочего дерева изолирует версии файлов, но не изолирует процессы и их сведения, сокеты и локальные сервисы, домашний каталог, общие кеши сборки, сокет Docker, временные каталоги, открытые порты, ресурсы. И главное: скрипты сборки. Постинсталляционный хук в npm или build-скрипт на Rust исполняется с правами агента, и ограничение путей в инструменте shell ему не указ. Границей может быть только механизм ядра.
\end{warnbox}

Политика ограничения задаёт: запись только в каталог санбокса и его временный каталог; чтение системы и тулчейна; запрет на остальной домашний каталог и локальные сервисы; запрет сигналов и сведений о чужих процессах; сеть закрыта, кроме локального egress-прокси, в котором живёт список разрешённых адресов. Агент никогда не получает shell хоста: каждая команда стартует внутри политики, и выйти из неё нельзя, не эксплуатируя ядро.

\subsection{Механизмы по платформам}

\begin{center}
\small
\begin{tabular}{p{2.3cm} p{5.3cm} p{6.0cm}}
\toprule
Платформа & Механизм & Что обеспечивает и чего нет \\
\midrule
macOS & Seatbelt (sandbox-exec), профиль на задачу; клон APFS; группа процессов со сторожем & файловая система, сеть, сигналы, сервисы: да; агрегатные лимиты памяти: нет, заменяет сторож; ядро общее \\
Linux, тысячи санбоксов & bubblewrap с user namespaces, seccomp, cgroups; overlay & всё, включая ресурсы; ядро общее \\
Linux, чужие агенты & gVisor & перехват системных вызовов, меньшая поверхность ядра \\
Linux, прочность VM & Firecracker, microVM на задачу & отдельное ядро; тысячи на хост штатно \\
macOS, прочность VM & macOS-гость в Virtualization & не больше двух гостей на хост; только для критических операций \\
\bottomrule
\end{tabular}
\end{center}

Seatbelt на macOS это тот же механизм, на котором стоит App Sandbox, и тот же, который используют собственные санбоксы вендоров агентских инструментов. Он даёт нулевую стоимость на процесс и тысячи одновременных профилей. Его ограничения честно перечислены в инвентаре (раздел 4.5): общее ядро и отсутствие агрегатных лимитов памяти.

\section{Слоты и очередь}

Санбокс дёшев: это каталог и политика, и в покое он не стоит ничего. Исполнение дорого: одновременно живых сборок на машине столько, сколько тянут процессор и память, для Xcode это десятки. Поэтому протокол различает санбоксы, которых могут быть тысячи, и слоты, число которых задаёт ёмкость среды. Назначение без свободного слота ждёт в очереди.

\begin{whybox}
Так устроен любой serverless: миллионы функций, сотни одновременных запусков. Полезный размер роя ограничен ёмкостью среды, а не числом контрибьюторов, и честнее сделать очередь частью протокола, чем обнаружить её в проде.
\end{whybox}

Отдельно про Xcode. DerivedData и кеш пакетов живут внутри санбокса, с прогретой базой, которая клонируется вместе с деревом. Симуляторы это разделяемый сервис, его нельзя честно посадить в санбокс; прогон на симуляторах и устройствах выполняется операцией раннера (раздел 4.4), а не из shell агента.

\section{Защищённые пути и операции раннера}

Тесты являются валидатором результата, поэтому защищаются сильнее кода. Тесты, конфигурация сборки, правила валидаторов и манифест проекта смонтированы во все санбоксы только на чтение. Их изменение проходит через задачи привилегированного класса (глава 6), санбоксы которых открываются с другой политикой.

Всё, что требует секретов проекта или общих сервисов, недоступно из shell вовсе. Подпись артефактов, публикация, прогоны на симуляторах и устройствах, доступ к данным предоставляются как узкие именованные операции доверенного раннера с очередью. Каждая операция авторизуется собственной записью журнала, имеет идентификатор и ключ идемпотентности и оставляет запись исхода: выполнена, отклонена или не установлен. Семантика повторов и поведение при смене аренды и остановке описаны в главе 7. Секреты проекта существуют только в раннере, секреты контрибьюторов не существуют на среде вовсе: это две разные половины инварианта I6. Агент вызывает операцию так же, как команду, но исполняет её не его санбокс, а раннер, у которого есть ключи и которому нельзя передать произвольную строку.

\begin{exbox}[Операции раннера в сценарии]
\texttt{sign\_artifact(hash)} подписывает сборку ключом проекта и возвращает хеш подписанного артефакта. \texttt{run\_device\_tests(candidate, suite)} ставит прогон в очередь на устройства и возвращает аттестацию. \texttt{publish(artifact, channel)} публикует. Ни в одной из них нет shell, и ни одна не принимает путь вне санбокса задачи.
\end{exbox}

\section{Receipts и инвентарь ограничений}

Каждая команда, которую агент отправляет в среду, получает идентификатор, digest самой команды и digest входного состояния санбокса, а также лимиты по времени, памяти, числу процессов и объёму вывода. По завершении среда выдаёт receipt: код выхода, дайджесты потоков вывода, длительность, потреблённые ресурсы, digest результирующего состояния. Receipts одной задачи связаны в хеш-цепочку, и голова цепочки входит в запись о ревизии (инвариант I10). Receipt удостоверяет происхождение, а не воспроизводимость и не полноту: он говорит, что среда получила такую-то команду и выдала такой-то результат, но ничего не говорит о том, что агент делал на своей машине вне протокола. Воспроизводимость требует отдельно сохранённых входов, digest окружения и тулчейна, зафиксированных зависимостей и, там, где они влияют на результат, внешних ответов; egress-прокси есть естественная точка, где внешние ответы фиксируются для повтора.

Модель команды и результата взята из API удалённого исполнения сборок: действие адресуется по содержимому, результат аттестован, повтор того же действия над тем же входом идемпотентен. Это позволяет верификатору повторить критическую команду и сравнить результат.

Инвентарь ограничений есть часть манифеста проекта: для каждого класса ограничения среда объявляет, обеспечивает она его полностью, частично или не обеспечивает. Задача объявляет требуемый класс. Несовпадение есть отказ, а не предупреждение (инвариант I9). На macOS инвентарь честно говорит: файловая система и сеть полностью, память частично, через сторож процесса.

\section{Светлый клиент журнала}

Среда исполняет только то, что подтверждено журналом (инвариант I1). Три точки входа:

\begin{itemize}
\item санбокс открывается по записи назначения с сертификатом кворума; каждая последующая команда несёт идентификатор задачи, подпись ноды и digest плана, и среда сверяет digest с записью;
\item операция раннера проводится по записи авторизации операции;
\item голова основной ветки продвигается по записи об интеграции с сертификатом кворума, атомарно, с проверкой, что текущая голова равна ожидаемой.
\end{itemize}

Среде не нужен весь журнал. Она хранит последний подписанный снапшот и проверяет включение конкретной записи по доказательству Меркла. Она хранит наивысшую принятую версию журнала как долговременное состояние безопасности и не принимает более раннюю (инвариант I8). И она проверяет актуальность разрешения на границе исполнения: запись, совершённая в прошлой эпохе или до остановки, не даёт права на исполнение сейчас, каким бы валидным ни был её сертификат. Компрометацию самой среды протокол не предотвращает; глава 9 разбирает, что именно это значит.

\section{Вход в санбокс}

Как человек с улицы попадает в санбокс на машине проекта и почему это не страшно. По шагам.

\begin{enumerate}
\item Идентичность без аккаунтов. У контрибьютора есть ключ ноды. Он делает join по адресу манифеста проекта, отправляет публичный ключ и объявленные оси, получает план и разницу полномочий. Кор-нода записывает ноду в журнал. Ни паролей, ни регистрации: запись о ноде есть, прав нет.
\item Права появляются только с назначением. Оркестратор создаёт запись назначения: ключ ноды, задача, digest плана, класс санбокса из контракта (платформа, процессор, память, диск, сеть, время) и слот. До совершения этой записи среда о контрибьюторе ничего не знает.
\item Санбокс заводится по записи, а не по запросу. Демон видит совершённое назначение через доказательство включения и провизионирует санбокс нужного класса с лимитами из контракта и egress-прокси. Санбокс живёт от назначения до закрытия задачи.
\item Аутентификация это сертификат на задачу. Среда держит свой CA. При совершении назначения демон выпускает короткоживущий SSH-сертификат для транспортного ключа ноды: principal равен идентификатору задачи, срок равен окну назначения, принудительная команда запускает демон по стандартным потокам ровно для этого санбокса. Сертификат приходит в ответ на назначение. Контрибьютор подключается своим ключом с этим сертификатом и оказывается сразу внутри санбокса. Отзыв не нужен: сертификат истекает сам, а закрытие задачи или остановка рвут сессии.
\item Каждая команда несёт digest плана, среда сверяет его с записью и выдаёт receipt.
\item Сетевой путь. Среда за NAT в офисе публикуется через маленький relay на VPS, который пробрасывает SSH в tailnet среды; среда на публичном адресе принимает только сертификаты, без паролей и без ключей в authorized\_keys.
\end{enumerate}

\begin{warnbox}[{SSH сразу в санбокс: безопасно при пяти условиях}]
Паттерн стандартный: так работают git-хостинги с принудительной командой, облачные среды разработки с SSH в контейнер, отладочные сессии CI. Небезопасным он становится, когда через SSH можно получить что-то, кроме санбокса. Условия. Первое: в сертификате опция restrict, никакого перенаправления портов, agent forwarding, X11, sftp, переменных окружения и pty; клиент не может запросить другую команду. Второе: команда исполняется от отдельного служебного пользователя без прав и без секретов; на Linux дальше user namespaces дают санбоксу свой uid, на macOS разделение процессов разных задач обеспечивает Seatbelt. Третье: только сертификаты от CA среды, без паролей и без ключей в файле; транспортный ключ сертифицирован ключом ноды и не совпадает с ключом подписи ревизий, один ключ на два протокола плохая практика. Четвёртое: поверхность sshd. OpenSSH работает от root, и уязвимость до аутентификации означает компрометацию машины; для публичного роя SSH-сервер встраивается в сам непривилегированный демон на библиотеке и принимает только exec протокола по stdio, root в приёме соединений не участвует; для L0 и корпоративного варианта достаточно OpenSSH с первым и третьим условиями. Пятое: границей всё равно остаётся санбокс. SSH это аутентификация, связывающая сессию с назначением, а не изоляция; если санбокс слаб, безопасный SSH не спасёт.
\end{warnbox}

Что защищает от человека с улицы, не аутентификация. Защищают граница ядра в санбоксе, отсутствие секретов на среде, allowlist сети, лимиты и слоты, и то, что в проект ничего не попадает без ревью другой ноды. Доказательства идентичности (глава 8) здесь необязательны: они меняют вес ноды и глубину ревью, а не право войти в свой санбокс.

\section{Shell, неотличимый от локального}

Удалённый shell есть единственный инструмент рабочей ноды, поэтому его качество определяет качество работы всего роя. Требование сформулировано так: для модели shell среды неотличим от локального, не в смысле обмана, а в смысле того, что замечать нечего.

Механика, которая это даёт, называется восстановленной сессией. Каждый вызов есть отдельный процесс: точный код выхода, раздельные потоки вывода, отмена по группе процессов без последствий для следующего вызова. При этом состояние сессии, каталог, переменные, функции, снимается после команды по отдельному дескриптору, мимо потоков вывода, и восстанавливается перед следующей командой. Модель делает \texttt{cd src} в одном вызове и оказывается там в следующем, как в терминале.

Мутирующие вызовы одной сессии сериализованы: две команды, одновременно меняющие каталог, не могут потерять изменение друг друга, они просто выполняются по очереди. Параллельность даётся именованными сессиями, как вкладками терминала. Длительные процессы не переживают вызов, поэтому они запускаются явно в фоне и живут по дескриптору. Вывод сверх лимита остаётся в санбоксе, модель получает начало и конец и дальше читает его тем же shell.

\begin{warnbox}
Восстановленная сессия есть механизм качества, а не безопасности. Команда модели может испортить снимок своей сессии; она не получает от этого ничего, чего у неё не было. Границей остаётся санбокс.
\end{warnbox}

Транспорт: SSH только для аутентификации по сертификату задачи (раздел 4.7), а поверх него собственный протокол по стандартным потокам ввода-вывода демона среды. Ни отдельной аутентификации, ни shell хоста, один тёплый канал с мультиплексированием вызовов. Демон среды принадлежит проекту, его хеш опубликован в манифесте, и клиент контрибьютора пинит ожидаемую идентичность. Тот же протокол поверх другого транспорта превращает среду в облаке или в контейнере в ту же самую среду с точки зрения агента.

\chapter{Журнал и инварианты}

Журнал есть плоскость управления в чистом виде: всё, что имеет силу в системе, существует как запись в нём. Эта глава описывает форму журнала, а затем разбирает десять инвариантов, потому что именно журнал делает их проверяемыми.

\section{Запись}

Запись содержит хеш предыдущей записи, тип, тело или хеши тел, и подписи. Тело сериализуется в канонический JSON: детерминированная форма, в которой одинаковые данные всегда дают одинаковые байты, а значит одинаковый хеш. Подписи Ed25519 ставятся поверх канонической формы; идентификатор ключа есть первые шестнадцать шестнадцатеричных символов хеша самого ключа. Набор пинов кор-ключей ротируется через перекрытие: сначала распространяется набор со старым и новым ключом, потом подписи переходят на новый, потом старый удаляется. Скомпрометированный ключ удаляется немедленно.

\begin{exbox}[{Запись назначения, набросок}]
\begin{lstlisting}
{
  "type": "assignment",
  "prev": "b3c1...",
  "task": "TASK-260901-7fk2",
  "node": "ed25519:...",
  "role": "producer",
  "plan_digest": "9a4e...",
  "lease": {"epoch": 41, "token": "f81d..."},
  "slot": "env-1/slot-07",
  "body_hash": "51aa...",
  "availability": ["replica-a", "replica-b"],
  "sig": [{"key_id": "0123456789abcdef", "signature": "..."}]
}
\end{lstlisting}
\end{exbox}

\section{Цепочка, дерево и снапшот}

Для записи $i$ с канонической формой $c(r_i)$ хеш вычисляется как $h_i = H(h_{i-1} \,\|\, c(r_i))$; у первой записи предшественник есть строка нулей. Над хешами записей строится дерево Меркла, и снапшот, подписанный кор-ключами, содержит корень дерева, длину журнала, голову и версию. История только дописывается: удаление и замена запрещены.

Что это даёт практически. Любой участник с последним снапшотом проверяет любую запись по доказательству включения. Изменение прошлого невозможно без пересчёта всех последующих хешей и новой подписи снапшота, то есть без кворума. Светлые клиенты, среда и контрибьюторы, не хранят журнал целиком.

\section{Типы записей}

\begin{center}
\small
\begin{tabular}{p{3.4cm} p{10.2cm}}
\toprule
Тип & Что фиксирует \\
\midrule
цель & область изменений, класс риска, бюджет; подпись кор-ноды \\
capability & часть цели, выданная оркестратору: пути, типы изменений, лимиты \\
контракт задачи & версия: write scope, зависимости, критерии, профиль свойств, оценка, класс, бюджеты; digest плана \\
зависимость & какие принятые артефакты задача использует и какие результаты предоставляет \\
назначение & задача, нода, роль, аренда, слот \\
ревизия & хеш коммита, база, голова цепочки receipts, подпись ноды \\
заморозка кандидатов & множество нод, из которых будет выбран ревьювер \\
случайность & значение маяка или VRF после заморозки и ревизии \\
кандидат & база $M$, неизменяемый кандидат $I$ \\
ревью кандидата & одобрение или находки, адресованные к $I$, версии контракта, политике, окружению \\
аттестация & свойство из профиля, исход (пройдена, провалена, не установлена), адрес, нода, верификатор при наличии \\
интеграция & пара $(M, I)$, ссылки на ревью и аттестацию \\
вестинг & закрепление начисления после $k$ эпох \\
аренда, эпоха & нода, эпоха, токен, предшественник, причина, чекпоинт, область \\
чекпоинт & хеш снимка сессии с сертификатом доступности \\
пауза, возобновление & экстренная остановка и снятие \\
состав & изменение множества кор-нод и ключей \\
операция раннера & авторизация именованной операции с идентификатором и ключом идемпотентности; отдельной записью исход \\
аттестация репутации & исход задачи для ноды \\
подтверждение идентичности & идентичность, класс доказательства, хеш свидетельства резолвера, окно; подписи валидаторов, каждый резолвил сам \\
сертификация ключа & ключ ноды, подписанный ключом идентичности, с началом и концом действия \\
\bottomrule
\end{tabular}
\end{center}

\section{Порядок, кворум и уровни}

Порядок записей задаёт секвенсер. Запись совершена, когда собран сертификат кворума уровня. На L0 секвенсер и единственный валидатор совпадают. На L1 секвенсер есть всегда включённый сервис проекта, которому доверяют в отношении порядка, а кор-ноды подписывают асинхронно; запись с $t$ подписями, $t > n/2$, совершена, и до этого момента среда её не исполняет. Порог большинства защищает от одиночного несанкционированного решения, но не от раздвоения истории при злонамеренном секвенсере (глава 3); L1 это признаёт, а не скрывает. На L2 порядок фиксируется византийски устойчивым консенсусом валидаторов, который даёт и то и другое при $f$ лжецах; среда исполняет только записи с $2f+1$ подписями.

\begin{whybox}[Почему секвенсер отдельно от валидаторов]
Консенсус между ноутбуками кор-команды непрактичен: живучесть требует, чтобы достаточное число нод было в сети постоянно. Секвенсер даёт живучесть, валидаторы дают безопасность. Так устроены роллапы, и так уже устроен реестр аудита, на котором журнал будет реализован.
\end{whybox}

Роль лидера консенсуса на L2 механическая: она ротируется по раундам, не требует модели и не совпадает с ролью оркестратора. Это важно повторять: оркестратор не лидер журнала, он один из отправителей записей.

\section{Доступность и материализация}

Тела задач, находки ревью, снимки сессий и артефакты хранятся адресуемо по содержимому, журнал содержит только хеши. Хеш включается в журнал только с сертификатом доступности: несколько независимых реплик подтвердили хранение объекта. Код тоже адресуется по содержимому: журнал хранит хеши коммитов, а git выступает хранилищем, в котором любая нода проверяет, что ревью, тесты и интеграция относятся к одному состоянию.

Журнал может материализоваться как борд задач в git-репозитории проекта: каждый переход есть коммит борда, сертификат кворума покрывает хеш коммита, полная история задачи есть история git. Журнал не дублирует борд, он сертифицирует его. Для команды это означает, что привычный борд остаётся и что ничего не нужно переносить: добавляются подписи и цепочка. То же относится к графу зависимостей задач и артефактов: он материализуется из записей о зависимостях и не является вторым источником истины.

\section{Десять инвариантов}

Каждый инвариант ниже сопровождён тем, что ломается без него. Это полезный тест при реализации: если предлагаемое упрощение делает поломку возможной, оно неприемлемо.

\begin{invbox}[I1. Эффект только по записи]
Ни одного эффекта на среде без совершённой записи авторизации, digest плана которой совпадает с digest плана вызова, и актуальность которой проверена на границе исполнения.

Без него: любой, кто дотянулся до среды, исполняет что угодно; оркестратор в одиночку открывает санбоксы; агент с расширенным планом работает под чужой задачей; операция, разрешённая в прошлой эпохе, исполняется после остановки.
\end{invbox}

\begin{invbox}[I2. Только сужение]
План задачи есть подмножество capability, capability есть часть цели; ревизия не изменяет ничего вне write scope задачи; условия приёмки меняются только новой версией контракта с отдельной авторизацией.

Без него: оркестратор создаёт задачи с полномочиями, которых цель не давала; отравленный субагент расширяет права родителя; задача трогает тесты, которые должна была только читать; критерий успеха ослабляется по ходу работы.
\end{invbox}

\begin{invbox}[I3. Интегрируется ровно проверенное]
Основная ветка продвигается только атомарной заменой $M$ на неизменяемого кандидата $I$, к которому адресованы финальное одобрение и аттестации, ссылающиеся на текущую версию контракта и политики.

Без него: две безопасные по отдельности задачи сливаются в небезопасную композицию, которую никто не проверял; одобрение, сделанное на старой базе, переезжает на кандидата над новой головой. Исполняемая модель в главе 11 показывает оба сценария как нарушения.
\end{invbox}

\begin{invbox}[I4. Суждение вне консенсуса]
Валидаторы проверяют только детерминированные предикаты над записями и хешами.

Без него: консенсус о правильности кода невозможен и превращается в политику; византийская устойчивость теряет смысл, потому что предмет голосования не верифицируем.
\end{invbox}

\begin{invbox}[I5. Независимость и случайность]
Продьюсер, ревьювер и исполнитель проверок независимы согласно объявленной политике проекта; отсутствие подходящих участников не ослабляет политику автоматически; случайность выбора фиксируется после заморозки кандидатов и хеша кандидата.

Без него: сговор продьюсера и ревьювера; одинаковые слепые зоны одной модели; подбор ревьювера оркестратором; тихое ослабление требований, когда некому ревьюить.
\end{invbox}

\begin{invbox}[I6. Ключи у владельца, секреты в раннере]
Ключи контрибьютора не покидают его машину, и подпись выполняется на машине агентского цикла; секреты проекта существуют только в раннере и недоступны ни из одного санбокса.

Без него: компрометация среды означает компрометацию всех контрибьюторов; ключи подписи проекта извлекаются через shell.
\end{invbox}

\begin{invbox}[{I7. Журнал истина, сессия кэш}]
Записи с устаревшей арендой отвергаются и сохраняются как расхождение.

Без него: два оркестратора одновременно; отравленный контекст переезжает вместе с ролью; смена оркестратора теряет состояние.
\end{invbox}

\begin{invbox}[{I8. Остановка дешёвая, откат невозможен}]
Любая кор-нода останавливает, кворум возобновляет; среда хранит наивысшую принятую версию журнала.

Без него: скомпрометированный оркестратор блокирует остановку; остановку обходят старым снапшотом.
\end{invbox}

\begin{invbox}[I9. Ограничение как факт]
Среда объявляет, что способна ограничить; задача, требующая полного ограничения, отказывает на частичном.

Без него: задача, рассчитанная на microVM, исполняется под Seatbelt, и никто об этом не знает.
\end{invbox}

\begin{invbox}[I10. Происхождение обращений]
Все обращения к среде через протокол и выданные ею результаты имеют проверяемое происхождение и связаны цепочкой receipts, голова которой входит в запись о ревизии. Протокол не удостоверяет полноту внутреннего контекста недоверенной рабочей ноды.

Без него: нельзя установить, что именно среда исполнила по чьему запросу; нельзя сравнить результат с верификатором. Вторая фраза не ослабление, а честная граница: во второй редакции инвариант обещал подотчётность всего, что модель видела, чего протокол наблюдать не может.
\end{invbox}

\chapter{Контракт задачи, полномочия и валидаторы}

Самая частая ошибка в системах с агентами: полномочия существуют неявно, как побочный эффект того, какие инструменты подключены к какому процессу, а условия приёмки живут в описании задачи, которое кто угодно может отредактировать. Здесь и то и другое существует только как записи, и это два разных объекта.

\term{Где это живёт сейчас.} Владелец этих определений есть спецификация конфигурации процесса: curator-playbook, \texttt{spec/process-configuration.md}, черновик v4.8. Проект объявляет свой процесс в \texttt{Playbook.json}: типы элементов и их workflow, охранные условия на переходах, наборы проверок, делегирование действий человека с потолками. Изменение процесса вступает в силу по подтверждению человека или в пределах прав, которые человек выдал заранее, и ядро исполняет только версию в силе. Детерминированные проверки на переходах дают охранные условия ядра, инварианты модуля доставки и статический валидатор \texttt{curator playbook check}. Эта глава остаётся объяснением протокольной формы: digest плана, контракт с версией, адрес аттестации и правила, которые кор-нода проверяет перед подписью. Где глава и спецификация расходятся, права спецификация.

\section{Два вопроса, два объекта}

\term{Digest плана} отвечает на вопрос, что агенту разрешено делать: привязки инструментов, гранты, лимиты. План компилируется из write scope и класса задачи, входит в запись назначения, и среда сверяет его на каждом вызове.

\term{Digest контракта} отвечает на вопрос, что агент должен получить и по каким правилам результат будет принят. Контракт задачи содержит родительскую цель, write scope, версии входных артефактов и зависимости, ожидаемые результаты и критерии приёмки, профиль подтверждаемых свойств, оценку, класс риска и бюджеты.

\begin{whybox}[Почему одного digest мало]
Во второй редакции у задачи был один digest, и он фиксировал полномочия. Ревью могло оказаться сделанным относительно одной формулировки критериев, а принятие относительно уже изменённой: оркестратор вправе создавать задачи, значит, вправе и переписать критерии по ходу, и ничто в журнале этого не ловило. Контракт с версией закрывает дыру: изменение критериев или профиля создаёт новую версию и требует авторизации того же уровня, что создание задачи. Обычная задача не может незаметно ослабить критерий собственного успеха, точно так же, как не может изменить тесты.
\end{whybox}

\section{Профиль свойств и исходы}

Профиль подтверждаемых свойств отвечает на четыре вопроса: какое свойство требуется подтвердить, каким способом, для какого артефакта и при каких допущениях. Для одной задачи это сборка, регрессионные проверки и независимое ревью. Для другой машинно проверяемое доказательство конкретного свойства плюс ревью того, что формальная формулировка соответствует исходной задаче. Проверки подтверждают разные свойства, и лестница из них не складывается: компиляция не заменяет проверку поведения, доказательство отсутствия переполнения не подтверждает правильность авторизации.

Исход каждой проверки есть одно из трёх: пройдена, провалена, не установлена. Таймаут или сбой инструмента даёт не установлена и никогда не превращается в принятие.

\begin{warnbox}
Реализация, в которой упавший или зависший инструмент проверки трактуется как отсутствие ошибок, есть самая тихая из возможных дыр. Три исхода вместо двух есть требование к каждому провайдеру проверки.
\end{warnbox}

\section{Адрес аттестации}

Финальное одобрение и каждая аттестация ссылаются на сочетание из четырёх частей: версия контракта, интеграционный кандидат $I$, версия политики проверки, digest окружения. Аттестация, ссылающаяся на устаревшую версию любой из четырёх, недействительна. Это защищает от ситуации, когда результат проверили относительно одной спецификации, а приняли относительно уже изменённой, и от ситуации, когда проверили одно состояние, а приняли другое.

\section{Иерархия}

\begin{center}
\begin{tikzpicture}[node distance=5mm]
\node[boxa,minimum width=62mm] (g) {Цель\\[1pt]\footnotesize подписана кор-нодой: область, класс риска, бюджет};
\node[boxa,minimum width=62mm,below=of g] (c) {Capability\\[1pt]\footnotesize выдана оркестратору: пути, типы изменений, лимиты};
\node[boxa,minimum width=62mm,below=of c] (t) {Задача\\[1pt]\footnotesize контракт (версия) и план (digest)};
\node[boxa,minimum width=62mm,below=of t] (r) {Ревизия и кандидат\\[1pt]\footnotesize диф внутри write scope; $I$ над головой $M$};
\draw[arr] (g) -- (c);
\draw[arr] (c) -- (t);
\draw[arr] (t) -- (r);
\node[lbl,right=10mm of c] {каждый уровень\\только сужает предыдущий};
\end{tikzpicture}
\end{center}

\term{Цель} создаётся и подписывается кор-нодой и есть единица человеческого решения. \term{Capability} выдаётся оркестратору как часть цели: допустимые пути, типы изменений, лимиты. \term{Задача} ссылается на родительскую цель и несёт контракт и план. \term{Ревизия} допустима, только если её диф лежит внутри write scope, а \term{кандидат} есть неизменяемое состояние, построенное из ревизии над текущей головой.

Полномочия задачи в журнале и полномочия агента на среде есть один объект с одним digest: запись назначения несёт digest плана, каждый вызов shell несёт digest плана, под которым работает агент, среда сверяет их и отвергает несовпадение. Оркестратор декомпозирует полномочия, но не создаёт их.

\begin{whybox}
В runtime агента план сессии уже есть: скомпилированный, неизменяемый в течение своей жизни, с digest, и каждый вызов инструмента уже несёт этот digest. Правило, что дочерняя сессия не может иметь больше, чем родительская, там тоже уже есть. Протокол лишь поднимает этот механизм на уровень задач.
\end{whybox}

\section{Классы риска}

Класс риска определяется детерминированно: по путям, манифестам зависимостей и типам операций. Привилегированный класс включает авторизацию и доступ, зависимости, конфигурацию сборки и CI, проверки, правила и манифест, подпись, публикацию, секреты, миграции данных. Такая задача требует двух независимых ревью либо одобрения человека из кор-команды, а её пути на среде доступны для записи только из санбокса задачи этого класса.

\begin{warnbox}[Остаточный риск]
Семантически вредная задача, лежащая в обычном классе и внутри своей области, остаётся возможной. Протокол ограничивает её областью, бюджетом, независимым ревью кандидата, вестингом и остановкой, но не устраняет. Это явное ограничение в модели угроз.
\end{warnbox}

\section{Write scope}

Две задачи с пересекающимися write scope не исполняются одновременно, и это правило валидатора. Область удерживается с момента назначения до интеграции или отмены задачи: следующая задача не начнёт менять те же пути, пока предыдущая ревизия ожидает принятия.

\begin{warnbox}[Что дизъюнктные области не обещают]
Дизъюнктные области предотвращают конкурирующую запись в одни и те же пути. Они не исключают семантических конфликтов: одна задача меняет контракт функции в библиотеке, другая добавляет её вызов в другом файле, и обе безопасны по отдельности. Такой конфликт ловит только проверка кандидата, построенного над текущей головой, и ничто раньше. Именно поэтому одобрение относится к кандидату, а не к ревизии (глава 7).
\end{warnbox}

Следствие для оркестратора: декомпозиция есть разбиение области цели на дизъюнктные части, а задачи, которые невозможно развести, исполняются последовательно.

\section{Зависимости и переиспользование}

Задача может использовать принятые артефакты других задач. Зависимость записывается в журнал: какие принятые артефакты задача использует и какие результаты предоставляет. Граф зависимостей задач и артефактов есть представление журнала, материализуемое из этих записей, а не второй источник истины. Индекс для поиска помогает найти полезный модуль, тест или прежний результат, но не определяет, принят ли он и применим ли к текущей версии контракта. При переиспользовании агент получает не просто полезный код, а его происхождение, контракт, зависимости и основания принятия.

\section{Правила валидаторов}

Кор-нода подписывает запись только после проверки предикатов. Все они детерминированы.

\begin{enumerate}
\item Переход состояния задачи допустим по автомату из главы 7.
\item Write scope лежит внутри capability, capability внутри цели; digest плана соответствует области и классу; версия контракта, на которую ссылается запись, есть текущая.
\item Продьюсер, ревьювер и исполнитель проверок различны согласно политике независимости; если политика невыполнима, записывается ожидание или эскалация, а не ослабление.
\item Множество кандидатов в ревьюверы заморожено записью до появления хеша кандидата, случайность получена после, ревьювер вычислен из идентификатора задачи, хеша кандидата и случайности.
\item Запись об интеграции содержит базу $M$, кандидата $I$, финальное одобрение и аттестации, адресованные ровно к $I$, текущей версии контракта, текущей политике и окружению, с исходом пройдена; текущая голова основной ветки равна $M$.
\item Назначение выдано действующей арендой с валидным fencing-токеном, в пределах лимитов capability и без пересечения write scope с удерживаемыми областями.
\item Каждый хеш тела имеет сертификат доступности.
\item Класс риска задачи и требуемый им процесс соблюдены.
\item Цель создана кор-нодой; предложения рабочих нод не исполняются, пока их не подпишет кор-нода.
\item Операция раннера имеет идентификатор и ключ идемпотентности, а повтор ссылается на тот же идентификатор.
\end{enumerate}

\begin{exbox}[Как проверяется правило 5]
Валидатор берёт из записи хеш $I$ и хеш $M$, читает голову основной ветки, сравнивает; берёт одобрение и аттестации, проверяет подписи и что каждая адресована к $I$ с той же версией контракта и политики; проверяет, что исход каждой проверки из профиля равен пройдена. Ноль обращений к модели.
\end{exbox}

\section{Чего валидаторы не проверяют}

Валидатор не проверяет качество кода и правильность ревью. Он проверяет, что ревью и проверки были, что они относятся ровно к тому состоянию и контракту, которые принимаются, и что их сделали независимые участники. Альтернатива, голосование о правильности кода, не имеет ни определения, ни механизма проверки. Протокол честно делит труд: машина проверяет процедуру, модели и люди судят о содержании, и первое делает второе подотчётным.

\chapter{Оркестратор и цикл задачи}

\section{Роль и её границы}

Оркестратор планирует: декомпозирует цели в задачи с контрактами, назначает продьюсеров, инициирует ревью и интеграцию. Это прикладная роль, а не роль в консенсусе. Все её действия суть записи, проходящие правила валидаторов. Полезно перечислить, чего оркестратор не может:

\begin{itemize}
\item не имеет доступа к среде помимо записей: у него нет shell, среда исполняет команды воркеров;
\item не выбирает ревьювера: выбор вычисляется из случайности, полученной после того, как подбирать уже нечего;
\item не интегрирует без одобрения и проверок, адресованных к кандидату;
\item не выходит за capability и не меняет условия приёмки иначе как новой версией контракта с авторизацией;
\item не может заблокировать остановку.
\end{itemize}

\section{Аренда и fencing}

Роль оформлена записью аренды: нода, эпоха, случайный fencing-токен, предшествующая аренда, причина (создание, плановая передача, принудительный захват, восстановление), чекпоинт как база передачи и область. Аренда продлевается heartbeat. Каждая запись оркестратора несёт эпоху и токен; валидаторы отвергают записи с устаревшей арендой, но сохраняют их как ветку расхождения для аудита (инвариант I7).

\begin{exbox}[Уснувший ноутбук]
Оркестратор на ноутбуке уходит в сон. Heartbeat пропадает, аренда истекает, роль по ротации переходит следующей кор-ноде с новой эпохой и новым токеном. Ноутбук просыпается и отправляет назначение с прежним токеном. Валидаторы его не подписывают, но записывают в ветку расхождений. Среда, даже если запись каким-то образом до неё дошла бы, не исполнит её без сертификата кворума и без проверки актуальности эпохи на границе исполнения.
\end{exbox}

\section{Смена эпохи}

Смена происходит по таймауту heartbeat, по исчерпанию ёмкости оркестратора, по решению кворума или по обнаружению нарушения правил. Новый оркестратор определяется детерминированной ротацией; византийская нода пропускается следующей сменой. Раз в $k$ эпох смена происходит холодно независимо от причины.

\term{Плановая передача} требует доказательства покоя: ввод заблокирован, нет незавершённых вызовов, дочерних процессов и незавершённых операций раннера, дайджест хранилища стабилен. После этого сессия передаётся преемнику, при необходимости с конвертацией между форматами агентов, и передача записывается в журнал с хешем снимка.

\term{Принудительный захват} есть отдельное явное действие, а не плановая передача, которая не удалась. Сессия не передаётся, предшественник отсекается токеном, преемник стартует из журнала.

\begin{whybox}[Почему при захвате сессия не передаётся]
Уходящий оркестратор может быть тем самым византийским узлом. Переданная им сессия может быть отравлена, и преемник, унаследовавший такой контекст, становится византийским сам: компрометация переезжает вместе с ролью и обесценивает арифметику $f$. Поэтому враждебная смена всегда холодная, а при плановой передаче переданная сессия недоверена и расхождение с журналом разрешается в пользу журнала.
\end{whybox}

\section{Состояние и чекпоинты}

Журнал содержит план, контракты, назначения и результаты, и этого достаточно для вступления преемника. Сессия модели есть кэш. Оркестратор периодически сохраняет снимок сессии в хранилище с сертификатом доступности и фиксирует его хеш записью чекпоинта; после сбоя преемник загружает снимок и докатывает записи журнала после него. Дисциплина для оркестратора: решения и обоснования пишутся в журнал структурно, по ходу работы; конвертация сессий между агентами теряет часть контекста, и только структурная запись делает сессию по-настоящему одноразовой. Эксперимент, проверяющий это напрямую, описан в главе 11: остановить оркестратора, удалить его сессию, проверить, продолжит ли преемник.

\section{Несколько оркестраторов}

Аренда несёт область, поэтому одновременно могут действовать несколько аренд с дизъюнктными областями. Задача, затрагивающая пересечение областей, принадлежит аренде, покрывающей объединение, либо создаётся кор-нодой напрямую.

Когда один борд процесса ведут несколько операторов, 28 сентября 2026 года решено направление: борд остаётся один на процесс, и у него или у совместной работы один \term{принимающий писатель} с fencing по эпохе. Каждая типизированная операция несёт актора, борд, версию конфигурации, область и эпоху аренды, ожидаемую ревизию и идентификатор операции. Эпоха проверяется при запуске агента, при финальной записи и при интеграции; писатель, потерявший полномочия, приостанавливает новые записи. Это та же аренда с fencing-токеном, применённая к одному борду. Решение идёт в трек waggle CM2 вместе с арендами областей (waggle, \texttt{spec/waggle.md} v5.2, §12), а назначается ли писатель по владению процессом или выбирается участниками, пока открыто. В терминах главного цикла узла оркестратор процесса есть его корневой оркестратор, а оркестраторы на других узлах работают арендованными поддеревьями и второго корневого не заводят (curator-keeper, \texttt{spec/keeper.md} v0.5.1, §3.2, §4).

\section{Цикл задачи}

\begin{center}
\begin{tikzpicture}[node distance=4mm,scale=0.9,every node/.style={transform shape}]
\node[box] (s1) {Авторизована};
\node[box,right=of s1] (s2) {Назначена};
\node[box,right=of s2] (s3) {Ревизия $R$};
\node[box,right=of s3] (s4) {Обратная связь};
\node[box,below=13mm of s1] (s5) {Кандидат $(M, I)$};
\node[box,right=of s5] (s6) {Ревью $I$};
\node[box,right=of s6] (s7) {Проверки $I$};
\node[box,right=of s7] (s8) {Интеграция};
\node[box,right=of s8] (s9) {Вестинг};
\draw[arr] (s1) -- (s2); \draw[arr] (s2) -- (s3); \draw[arr] (s3) -- (s4);
\draw[arr] (s4.south) -- ++(0,-5mm) -| (s5.north);
\draw[arr] (s5) -- (s6); \draw[arr] (s6) -- (s7); \draw[arr] (s7) -- (s8); \draw[arr] (s8) -- (s9);
\draw[arr] (s4.north) -- ++(0,5mm) -| node[lbl,above,pos=0.25] {не более $N$ раундов} (s3.north);
\draw[arr] (s7.south) -- ++(0,-4mm) -| node[lbl,below,pos=0.3] {отклонено или голова ушла вперёд: пересборка кандидата} (s5.south);
\end{tikzpicture}
\end{center}

\subsection{Авторизация и назначение}

Оркестратор декомпозирует цель в задачи с контрактами; задача совершается записью с сертификатом кворума. Рабочие ноды сообщают в heartbeat остаток ёмкости и окно его сброса. Оркестратор назначает продьюсера с учётом ёмкости, свободного слота и удерживаемых областей. Среда открывает санбокс по записи назначения.

\subsection{Ревизия и обратная связь}

Продьюсер работает в санбоксе и представляет хеш коммита $R$ над базой $B$. Подпись ревизии ставится на машине продьюсера, ключ никогда не покидает ноду (инвариант I6). Голова цепочки receipts входит в запись.

Ранняя обратная связь ревьювера на $R$ допустима и полезна: продьюсер получает находки и представляет новую ревизию, число оплачиваемых раундов ограничено $N$. Но обратная связь не есть одобрение. Одобрение относится только к кандидату.

\subsection{Интеграционный кандидат}

\begin{center}
\begin{tikzpicture}[node distance=9mm]
\node[box] (b) {$B$\\[1pt]\footnotesize база задачи};
\node[box,right=of b] (r) {$R$\\[1pt]\footnotesize ревизия продьюсера};
\node[box,below=8mm of b] (m) {$M$\\[1pt]\footnotesize текущая голова main};
\node[box,right=of m] (i) {$I$\\[1pt]\footnotesize кандидат = слияние $M$ и $R$};
\draw[arr] (b) -- (r);
\draw[arr] (b) -- node[lbl,left] {main ушла вперёд} (m);
\draw[arr] (m) -- (i);
\draw[arr] (r) -- (i);
\node[lbl,right=4mm of i,align=left,text width=52mm] {одобрение и проверки адресованы к $I$, версии контракта, политике и окружению; после сдвига головы кандидат перестраивается, ревью и проверки повторяются};
\end{tikzpicture}
\end{center}

Среда строит неизменяемого кандидата $I$ над текущей головой $M$ и записывает пару $(M, I)$. Множество кандидатов в ревьюверы замораживается до появления хеша $I$; случайность получается после, из порогового маяка кор-нод или проверяемой случайной функции; ревьювер вычисляется из идентификатора задачи, хеша $I$ и случайности, и любой валидатор пересчитывает выбор.

\begin{whybox}[Почему одобрение нельзя переносить с ревизии на кандидата]
Во второй редакции одобрение $R$ переносилось на $I$, если совпадал patch-id дифа. Это было ошибкой. Одинаковый патч на другой базе может иметь другой смысл: одна задача изменила контракт функции, другая добавила её вызов, патч второй не изменился, но предположение, на котором стояло ревью, уже неверно. Исполняемая модель фрагмента интеграции (глава 11) находит ровно этот сценарий как нарушение, когда правило привязки к кандидату отключено. Поэтому финальное одобрение относится к неизменяемому $I$. Для людей повторное ревью после каждого перебазирования было бы неподъёмным, и очереди интеграции его не делают; для моделей оно стоит токены, и правило посильно. Перенос допустим только как расширение вне гарантии, при детерминированном условии: дельта базы $B..M$ не пересекается ни с write scope задачи, ни с множеством путей, которые продьюсер читал по receipts. Проверки на $I$ повторяются всегда.
\end{whybox}

\subsection{Ревью и проверки кандидата}

Ревьювер получает диф $M..I$ и подписывает одобрение, адресованное к $I$, версии контракта, политике и окружению. Проверки из профиля контракта выполняются на $I$ по инициативе ноды, третьей по отношению к продьюсеру и ревьюверу; критические наборы повторяет верификатор, если он есть. Каждая аттестация имеет один из трёх исходов; принятие требует исхода пройдена по всем свойствам профиля. При отклонении кандидата продьюсер получает находки, и цикл возвращается к ревизии.

\subsection{Интеграция}

Оркестратор вносит запись об интеграции с парой $(M, I)$. Валидаторы проверяют правило 5. Среда, получив сертификат кворума, атомарно заменяет голову основной ветки с $M$ на $I$. Если голова уже не равна $M$, сертификат неприменим: кандидат перестраивается, ревью и проверки повторяются (инвариант I3). Порядок интеграции задаёт очередь.

\subsection{Вестинг, превышение $N$, замена}

Начисление за задачу закрепляется, когда $I$ прожил $k$ эпох без отката. При превышении $N$ раундов оркестратор обязан выбрать из фиксированного меню: передекомпозировать задачу, сменить продьюсера, сменить ревьювера или эскалировать человеку; выбор записывается. Замена продьюсера происходит при потере ноды, исчерпании ёмкости или по решению из меню; преемник получает сессию предшественника с правилом недоверия к ней.

\section{Внешние эффекты}

Операции раннера, подпись, публикация, устройства, данные, имеют эффекты вне среды, и журнал авторизаций сам по себе не определяет их семантику при сбоях. Сценарий: публикация авторизована, раннер её выполняет и падает до сохранения результата; новый оркестратор видит авторизацию без завершения. Можно ли повторить?

\begin{itemize}
\item Каждая операция имеет идентификатор и ключ идемпотентности. Повтор ссылается на тот же идентификатор и параметры; для получателей, поддерживающих идемпотентность, повтор не создаёт дополнительного эффекта.
\item Исход записывается как выполнена, отклонена или не установлен. Сбой раннера до записи исхода даёт не установлен.
\item При исходе не установлен слепой повтор запрещён: требуется сверка фактического состояния получателя либо разрешение человеком, и результат сверки записывается.
\item Актуальность разрешения проверяется на границе исполнения. Операция, авторизованная в прошлой эпохе или до остановки, не исполняется без новой авторизации. Уже запущенная операция при смене аренды доводится до записи исхода; при остановке прерывается там, где это безопасно, иначе доводится и записывается.
\end{itemize}

\begin{warnbox}
Старая корректная авторизация не равна действующему разрешению на любое будущее исполнение. Правило актуальности проверяется демоном среды и раннером в момент исполнения, а не только валидаторами в момент записи.
\end{warnbox}

\section{Одна задача, все записи}

\begin{exbox}[Задача из сценария, от цели до вестинга]
\begin{enumerate}
\item Цель: перенос модуля синхронизации; область \texttt{Sources/Sync/**}; класс обычный. Подпись кор-ноды, QC.
\item Capability оркестратору: та же область, не больше шести одновременных задач.
\item Задача \texttt{TASK-7fk2}: контракт версии 1 (write scope \texttt{Sources/Sync/Transport/**}, критерии, профиль: сборка, регрессия, ревью), оценка 5, digest плана.
\item Назначение: нода контрибьютора А, роль продьюсера, аренда эпохи 41, слот 7. Среда открывает санбокс.
\item Ревизия $R$ над базой $B$, подпись А, голова цепочки receipts.
\item Обратная связь: нода В, две находки. Ревизия $R'$ от А.
\item Кандидат: голова $M$, слияние $I$; записана пара $(M, I)$.
\item Заморозка кандидатов в ревьюверы: ноды Б, В, Г. Случайность эпохи 41 после хеша $I$.
\item Ревью $I$: нода В, одобрено; адресовано к $I$, контракту версии 1, политике, окружению.
\item Аттестации: нода Г инициировала сборку и регрессию на $I$, исход пройдена; верификатор подтвердил критический набор.
\item Интеграция: пара $(M, I)$, QC. Среда: голова была $M$, стала $I$.
\item Вестинг через $k$ эпох; аттестации репутации А, В, Г.
\end{enumerate}
Двенадцать записей, из них ни одной, которую можно было бы подделать одной подписью на L1 и выше, и ни одной, для которой валидатору понадобилась бы модель.
\end{exbox}

\chapter{Ноды, независимость, остановка}

\section{Две формы клиента}

Рабочая нода это агентский цикл контрибьютора. Протокол допускает две формы клиента, и выбор между ними определяется тем, кому доверяет проект.

\term{Публичный рой} использует минимальный runtime, план которого содержит ровно две привязки: удалённый shell среды и подпись ревизий. Runtime не поставляет инструментов сам по себе, поэтому донорский профиль есть свойство конструкции, а не настройка: расширить его нельзя, не изменив digest плана, а digest плана сверяет среда на каждом вызове.

\term{Доверенная организация}, компания или кооператив, может вместо этого подключить среду как сервер инструментов внутри родного харнеса вендора. Донорский профиль там держится политикой организации: у агента остаются локальные инструменты, и проект доверяет тому, что организация их контролирует.

\begin{whybox}[Почему для публичного роя не подходит харнес вендора]
Харнес вендора всегда несёт локальный shell и локальные файлы. Задача от незнакомого проекта есть недоверенный ввод, и инжект в ней получает в распоряжение машину контрибьютора. Ограничить харнес снаружи нельзя надёжно; можно только не поставлять лишнего. Минимальный runtime с двумя привязками решает это по построению. Цена: подписки на модели можно использовать только там, где условия вендора допускают сторонний runtime; иначе источником ёмкости становятся ключи API и локальные модели.
\end{whybox}

\term{Решение 28 сентября 2026 года.} Собственный агентский цикл отложен: и оркестраторы, и рабочие ноды работают на существующих харнесах (Claude Code, Codex, Muse, харнес над локальной моделью), и независимость от харнеса есть принцип дизайна (curator-keeper, \texttt{spec/keeper.md} v0.5.1, §7.7). Минимальный runtime откладывается вместе с циклом. Пока его нет, донорский профиль публичной формы держится переопределением инструментов существующего харнеса; именно этот путь спецификация главного цикла называет возможным первым путём к ограниченным рабочим агентам. Это слабее свойства конструкции, потому что харнес по-прежнему несёт локальный shell и файлы, а профиль держится ровно настолько, насколько надёжно переопределение. Врезка выше остаётся верной как довод о том, чего переопределение должно достичь.

Обе формы говорят с проектом одним протоколом ноды: присоединение по ключу, heartbeat, назначение, сессии удалённого исполнения, представление ревизии, ревью, выход.

\section{Ёмкость любого происхождения}

Ресурс роя есть inference-ёмкость, а не подписка. Источники: подписки, где условия это допускают; ключи API; локальные модели на собственном железе контрибьютора. Разнообразие происхождения не только расширяет предложение, но и снижает корреляцию ошибок ревью: два ревьювера на разных семействах моделей пропускают разное.

\begin{whybox}[Что на самом деле отдаёт контрибьютор]
Folding@home работает потому, что простаивающий процессор ничего не стоит владельцу. Квота подписки не простаивает, её хотят потратить сами. Простаивает только то, что сгорит: неиспользованный остаток окна квоты перед сбросом. Это и есть главный ресурс публичного роя. Heartbeat сообщает остаток, окно сброса и класс стоимости; планировщик берёт истекающую ёмкость первой. Нода ограничивает расход на проект за окно, чтобы рой не съел то, что владелец хотел оставить себе.
\end{whybox}

\section{Подпись на машине владельца}

Ключ ноды никогда не оказывается на среде. Среда возвращает хеши состояний и вывода; провайдер на машине ноды подписывает ревизию и receipts (инвариант I6). В runtime это устроено так: ядро, которое обрабатывает вывод модели, не имеет доступа к ключам вовсе; ими владеет брокер, и подпись есть операция брокера, а не инструмента модели. Модель просит подписать хеш, брокер подписывает то, что видит в журнале задачи, и не подписывает ничего другого. Собственный runtime отложен (раздел 8.1), поэтому на существующих харнесах роль брокера должен играть компонент на машине владельца вне процесса модели. Ближе всего к этому хост сессий: по waggle он уже держит ключ сессии и подписывает за агента, который ключа не видит (waggle, \texttt{spec/waggle.md} v5.2, §4.2).

\begin{center}
\begin{tikzpicture}[node distance=8mm]
\node[boxa,text width=34mm] (core) {Ядро агента\\[1pt]\footnotesize обрабатывает вывод модели, ключей нет};
\node[boxg,text width=34mm,right=17mm of core] (broker) {Брокер\\[1pt]\footnotesize держит ключи, подписывает хеши};
\node[boxw,text width=34mm,right=17mm of broker] (env) {Среда\\[1pt]\footnotesize возвращает хеши, ключей не видит};
\draw[arr] (core) -- node[lbl,above] {подпиши\\хеш $R$} (broker);
\draw[darr] (broker) -- node[lbl,above] {команды,\\хеши} (env);
\end{tikzpicture}
\end{center}

\section{Донорский профиль и инжект}

Задача от проекта есть недоверенный ввод для агента. В публичной форме клиента план из двух привязок означает, что инжект в задаче не имеет локальных инструментов, к которым мог бы обратиться: он может только исполнять команды в санбоксе среды, куда и так был направлен, и просить подписать хеши, которые брокер сверяет с журналом. Бюджеты на проект за окно ограничивают ущерб от задачи, которая просто жжёт ёмкость.

\section{Конфиденциальность}

Рабочая нода видит содержимое своего санбокса и отправляет его провайдеру своей модели. Протокол не скрывает код от контрибьюторов и их провайдеров, и это нужно понимать до присоединения. Манифест проекта объявляет допустимых провайдеров и условия хранения; проект, которому это неприемлемо, ограничивает контрибьюторов или провайдеров, и это его политика, а не свойство протокола.

\section{Три оси независимости}

Два владельца на одной модели с одним харнесом ошибаются одинаково. Поэтому независимость определяется по трём осям:

\begin{itemize}
\item административный домен, против сговора и множественных личностей;
\item семейство модели, против общих слепых зон;
\item тулчейн, против одинаковых ошибок среды и анализаторов.
\end{itemize}

Независимость есть политика проекта с проверяемыми основаниями, а не пожелание. Проект объявляет, какие оси обязательны, и для каждого сведения о ноде говорит, к какому из трёх видов оно относится: подтверждается (ключ ноды и способ оплаты ёмкости подтверждают владельца в меру доверия к ним), принимается на доверии (объявленное семейство модели) или служит эвристикой. Различие строк в объявлении не доказывает различия владельцев.

\begin{warnbox}[Политика не ослабляется сама]
Если подходящих участников нет, задача ждёт, эскалируется либо получает отдельно авторизованное изменение условий, записанное в журнал. Реализация, которая молча снижает требования, когда некому ревьюить, ломает инвариант I5 в самом уязвимом месте.
\end{warnbox}

Отдельно различаются независимость ревью и независимость исполнения проверок. Третья нода, запустившая проверки в той же среде, даёт разделение ролей, но не независимое подтверждение честности среды; его даёт только верификатор на другой машине. Формула вероятности пропуска в главе 9 предполагает независимость ошибок; три оси делают ошибки менее коррелированными, но не независимыми, и это надо помнить при чтении любых чисел.

\section{Идентичность и привязка ключей}

Ось владельца из объявленной становится проверяемой через доказательства идентичности. Это не отдельная подсистема, а способ заполнить графу подтверждается в политике независимости.

\subsection{Двустороннее доказательство}

Владелец ноды имеет долгоживущий \term{ключ идентичности} и доказывает его принадлежность внешней сущности с двух сторон: заявлением я есть X у провайдера P, подписанным ключом, и артефактом на стороне провайдера, в котором лежит отпечаток ключа. Провайдеры суть \term{резолверы} с одним интерфейсом: по строке идентичности вернуть набор отпечатков ключей и хеш свидетельства.

\begin{likbez}[{Прецеденты: Keybase, DNS, well-known}]
Схема с заявлением у себя и отпечатком у провайдера известна по Keybase, где пользователь доказывал владение аккаунтами соцсетей, доменом и репозиторием. Публикация ключей в DNS стандартизована в DANE и в записях OPENPGPKEY; Bluesky привязывает handle к идентичности TXT-записью домена; Nostr публикует ключи по well-known адресу. Ничего из этого не нужно изобретать, нужно выбрать набор резолверов и определить, как результат их работы попадает в журнал.
\end{likbez}

Стартовый набор резолверов: TXT-запись домена под DNSSEC для организаций; well-known по HTTPS; ключи аккаунта на хостинге кода, потому что контрибьюторы это разработчики и ключ у них уже опубликован. Публикации в социальных платформах как необязательный класс с меньшим весом. Системы имён с блокчейном как плагины позже, если понадобятся. Протоколу важен только результат резолвера и его свидетельство.

\subsection{Подтверждение как запись журнала}

Резолв есть внешний вызов и недетерминирован между валидаторами: DNS не разошёлся, лимит API, сеть. Поэтому идентичность считается подтверждённой не потому, что кто-то проверил, а потому, что $t$ валидаторов независимо получили и записали одинаковый хеш свидетельства в одном окне. Расхождение даёт не подтверждено, без исключений. Это то же разделение подписанного свидетельства и независимо установленного факта, что и в главе 9, и оно ложится в секвенсер с асинхронными подписями без нового механизма.

Свидетельство хранится в журнале с датой. Пост удалили, API закрыли, домен временно недоступен: зафиксированное наблюдение остаётся, а политика проекта задаёт каданс переподтверждения. Первое увиденное свидетельство пинится: смена ключа в доказательстве есть событие с записью, а не молчаливое обновление.

\subsection{Ключ идентичности и ключи нод}

Ключ идентичности долгоживущий, он лежит в доказательствах и меняется редко. Ключи нод и устройств короткоживущие и сертифицируются ключом идентичности записью в журнале с началом и концом действия. Нода ротирует ключи, не переделывая доказательств; репутация привязана к идентичности и переживает ключи и проекты. Ротация через перекрытие и отзыв удалением у нас уже есть в реестре Curator для кор-ключей, здесь тот же набор пинов уровнем выше. Отзыв прост: убрать отпечаток из TXT-записи или well-known, и первый валидатор, не нашедший его, записывает расхождение.

\subsection{Подпись коммитов и сообщений}

Ревизии подписываются ключом ноды в формате SSH-подписи git. Файл допустимых подписантов для среды выводится из журнала: ключ допустим, если сертифицирован подтверждённой идентичностью и назначен на задачу. Хостинг кода покажет подпись как проверенную побочно, но проверяющим является журнал, а не платформа. Сообщения протокола уже подписаны ключами нод; добавляется только цепочка до ключа идентичности.

\subsection{Что это покупает и чего нет}

Покупает: проверяемую независимость владельцев (два домена или две сертифицированные идентичности, а не две строки в объявлении); рост цены множественных личностей, потому что домен или старый аккаунт дороже ключа за секунду; переносимую репутацию; отзыв; подотчётность людям, а не отпечаткам. Симметрично для проектов: ключи кор-команды публикуются под доменом организации, и контрибьютор при присоединении проверяет подписантов манифеста через домен.

\begin{warnbox}[Чего доказательство не удостоверяет]
Контроль над доменом не есть честность, аккаунт можно купить, DNS без DNSSEC угоняется. Доказательство удостоверяет контроль, а не поведение. Поэтому политика назначает вес по классу: домен под DNSSEC выше социального поста, а псевдонимная нода допускается с меньшим весом и большей глубиной ревью, но не запрещена. Запрет псевдонимов отрезал бы часть предложения ёмкости, не купив безопасности.
\end{warnbox}

В план MVP слой идентичности не входит: журнал уже умеет хранить ключи и доказательства как записи, и резолверы добавляются, не трогая инвариантов. Это первый релиз после MVP.

\section{Репутация}

Нода идентифицируется ключом и объявляет три оси. Стоимость создания ноды равна стоимости ёмкости, что даёт естественный барьер против множественных личностей, но нода на дешёвой модели тоже нода, поэтому барьер дополняется репутацией.

После каждой задачи проект выдаёт ноде подписанную аттестацию: принято, отклонено, вердикт ревью совпал с итогом или нет. Репутация определяет глубину проверки, как адаптивная репликация в BOINC: новой ноде назначаются два или три независимых ревьювера, ноде с историей один. Она определяет и вес при выборе в ревьюверы. Пропорциональное вознаграждение при полной устойчивости к множественным личностям недостижимо, поэтому репутация снижает риск, а не устраняет его.

Аттестации переносимы между проектами и проверяются офлайн. Проект может выдавать их собственным подставным нодам, поэтому чужие аттестации взвешиваются репутацией выдавшего проекта.

\section{Остановка}

Любая кор-нода выпускает сертификат экстренной остановки: подпись одного ключа, эпоха и монотонный номер против повтора. Среда по нему немедленно переходит в безопасный режим: не открывает санбоксы, не проводит операции раннера, не продвигает основную ветку. Позднее остановка упорядочивается журналом. Возобновление требует сертификата кворума уровня.

\begin{whybox}[{Почему одна подпись останавливает, а кворум возобновляет}]
Асимметрия намеренная: остановка дешёвая и безопасная, продолжение дорогое. В системе, где порча основной ветки хуже простоя, целостность важнее доступности. Нода, злоупотребляющая остановкой, исключается записью о составе, которая сама проходит кворум. Скомпрометированный оркестратор не может заблокировать остановку, потому что она не проходит через него, а откат журнала на версию до остановки невозможен по инварианту I8.
\end{whybox}

\chapter{Модель угроз}

Модель угроз есть договор между протоколом и его пользователем: при каких допущениях что гарантируется. Всё, что не перечислено как гарантия, не гарантируется, и это сказано явно.

\section{Участники и допущения}

\begin{center}
\small
\begin{tabular}{p{3.2cm} p{10.4cm}}
\toprule
Участник & Допущение \\
\midrule
Рабочая нода & не доверена; может быть злонамеренной, скомпрометированной или отравленной инжектом через содержимое задачи, репозитория или зависимостей \\
Кор-нода, оркестратор & L0: доверены; L1: не более меньшинства скомпрометированных ключей при честном секвенсере; L2: не более $f$ злонамеренных из $3f+1$ \\
Среда & доверена в отношении целостности на всех уровнях, если не аттестована аппаратно; протокол защищает от кор-нод и агентов, но не от владельца среды \\
Провайдер модели & видит содержимое санбокса; доверен в отношении конфиденциальности в пределах манифеста \\
Сеть & не доверена; сообщения подписаны, повтор отсекается эпохой и номерами \\
\bottomrule
\end{tabular}
\end{center}

\section{Что гарантируется}

При соблюдении допущений уровня:

\begin{itemize}
\item ни одного эффекта на среде без совершённой и актуальной авторизации (I1);
\item ни одного продвижения основной ветки без финального одобрения и аттестаций, адресованных ровно к продвигаемому кандидату и текущему контракту (I3);
\item полномочия и условия приёмки задачи не превышают и не ослабляют выданного (I2);
\item оркестратор не выбирает ревьювера и не интегрирует в одиночку (I4, I5);
\item остановка доступна одной честной кор-ноде (I8);
\item каждое обращение к среде через протокол имеет проверяемое происхождение (I10).
\end{itemize}

\section{Чего не гарантируется}

\begin{warnbox}[Явные ограничения]
\begin{itemize}
\item Правильность суждения модели. Ревью независимо и подотчётно, но не доказано верным.
\item Предотвращение компрометации среды. Она обнаруживается по расхождению хешей и receipts, а при наличии верификатора по расхождению тестов, но не предотвращается.
\item Конфиденциальность кода от контрибьюторов и их провайдеров.
\item Устойчивость к побегу из санбокса через уязвимость ядра. Санбоксы не разделяют записываемых каталогов, а всё ценное вынесено за операции раннера, чтобы побег давал как можно меньше.
\item Единственность истории на L1 при злонамеренном секвенсере.
\item Полнота наблюдения за тем, что недоверенная нода делает вне протокола.
\item Живучесть при злонамеренном большинстве кор-нод.
\item Справедливость экономики.
\end{itemize}
\end{warnbox}

\section{Сценарий: злонамеренный оркестратор}

Оркестратор получил роль и раздаёт вредные задачи или занимается вандализмом. Четыре слоя:

\begin{enumerate}
\item У него нет рук. Среда исполняет только команды рабочих нод по назначениям с сертификатом кворума. Сам оркестратор действует лишь записями.
\item Каждая запись проходит валидаторов. Он не выбирает ревьювера, не интегрирует без чужого ревью и аттестации тестов с третьей ноды, не превышает лимиты, не выходит за capability, не трогает через воркеров защищённые пути.
\item Вредная по смыслу задача, прошедшая правила, исполняется и ревьюится посторонними нодами, а интеграцию подписывают другие кор-ноды, которые видят поток записей. Любая из них ставит систему на паузу одной подписью, после чего кворум сменяет оркестратора и исключает его из состава.
\item Ущерб обратим. Хеши в журнале и git позволяют откатить состояние к последнему вестированному, начисления вандала не вестятся, его аттестации отрицательны.
\end{enumerate}

Остаточный риск: правильно исполненная неправильная задача, интегрированная до того, как её заметили. Он ограничен областью capability, бюджетом, вестингом и остановкой, но не устранён.

\section{Сценарий: злонамеренный воркер}

Воркер получил задачу и пытается навредить. Слои: санбокс ограничивает эффекты областью задачи, без сети и без секретов; правило write scope отвергает ревизию, вышедшую за область; случайный ревьювер другого домена на другой модели видит диф; тесты на кандидате гоняет третья нода; вестинг задерживает начисление; аттестация репутации снижает вес ноды и увеличивает глубину проверки её следующих задач. Что воркер может сделать наверняка: сжечь свою ёмкость впустую и получить отрицательную аттестацию.

\section{Сценарий: скомпрометированная среда}

Среда может показать ревьюверу один диф, а интегрировать другой; вернуть поддельный результат проверок; выполнить другую команду; утечь код. Протокол это не предотвращает. Важно и то, чего он не обещает: receipt есть подписанное свидетельство самой среды, а не независимо установленный факт, и если среда скомпрометирована, её подпись не доказывает правдивость отчёта. Обнаружение поэтому условно: оно требует независимого наблюдения, верификатора на другой машине, который повторяет критические проверки, или собственных git-объектов у ревьювера, по которым можно сверить диф с тем, что записано. Журнал пинит хеши коммитов, одобрений и аттестаций, receipts несут дайджесты команд и состояний, и это делает расхождение выявимым тем, у кого есть с чем сравнить. Аппаратная аттестация исполнения могла бы перевести это из обнаружимого в предотвращаемое, но на macOS её нет, и первое развёртывание честно живёт с доверенной средой.

\section{Числа}

Пусть доля злонамеренных рабочих нод равна $p$, честный ревьювер обнаруживает злонамеренную ревизию с вероятностью $d$, ревьюверов $k$. Злонамеренная ревизия проходит ревью, если каждый ревьювер либо злонамерен, либо честен и пропустил:
\[
q = p + (1 - p)(1 - d), \qquad P = q^{k}.
\]

\begin{center}
\begin{tabular}{cc ccc}
\toprule
$p$ & $d$ & $k=1$ & $k=2$ & $k=3$ \\
\midrule
0.1 & 0.7 & 0.370 & 0.137 & 0.051 \\
0.1 & 0.9 & 0.190 & 0.036 & 0.007 \\
0.3 & 0.7 & 0.510 & 0.260 & 0.133 \\
0.3 & 0.9 & 0.370 & 0.137 & 0.051 \\
\bottomrule
\end{tabular}
\end{center}

Два вывода. Первый: вероятность обнаружения $d$ влияет сильнее, чем доля злонамеренных $p$, и $d$ есть свойство модели и протокола ревью, которое протокол не улучшает, а умножает: независимыми ревьюверами, тестами, верификатором, вестингом и репутацией. Второй: формула предполагает независимость ошибок, а три оси делают ошибки лишь менее коррелированными. Таблица есть иллюстрация, а не оценка безопасности; реальная величина измеряется на данных проекта.

\section{Обнаруживаемое и предотвращаемое}

\begin{center}
\small
\begin{tabular}{p{6.6cm} p{6.6cm}}
\toprule
Предотвращается & Обнаруживается, но не предотвращается \\
\midrule
эффекты без авторизации & компрометация среды \\
интеграция непроверенного состояния & семантически вредная задача внутри своей области \\
превышение полномочий & пропуск дефекта ревьювером \\
выбор ревьювера оркестратором & \\
повтор и откат записей & \\
блокировка остановки & \\
\bottomrule
\end{tabular}
\end{center}

Правый столбец есть честная цена того, что суждение вынесено за консенсус. Левый столбец есть то, чего не даёт ни одна из существующих систем параллельных агентов.

\part{Как это построить}
\chapter{Существующие компоненты}

\begin{warnbox}[Глава заменена роадмапом экосистемы]
Состояние на 29 сентября 2026 года. Эта глава заменена роадмапом экосистемы (вики, \texttt{roadmap/ecosystem-roadmap.ru.md}, черновик v4: порядок приоритетов, модуль хоста сессий SH, этапы M0--M8, треки PC, KP и CM) и репозиториями модулей; карта репозиториев по плоскостям дана в \texttt{README.md} вики. С тех пор, как глава написана, хостинг сессий выносится в модуль agent-session-host, процесс описывает конфигурация процесса (curator-playbook), доверие вынесено в curator-trust, главный цикл узла в curator-keeper, сообщения и аренды между оркестраторами в waggle, а собственный агентский цикл и минимальный runtime на нём отложены (curator-keeper, §7.7). Текст главы сохранён как история; где он расходится с роадмапом или спецификацией модуля, правы они.
\end{warnbox}

Протокол не строится с нуля. Четыре компонента уже существуют, и большая часть механики протокола есть их композиция. Эта глава описывает каждый в той мере, в какой он нужен для протокола, и даёт карту соответствий.

\section{Карта соответствий}

\begin{center}
\small
\begin{tabular}{p{4.2cm} p{4.4cm} p{4.6cm}}
\toprule
Понятие протокола & Компонент & Что берётся \\
\midrule
запись, подпись, каноническая форма & реестр аудита Curator & CCJ-1, конверт Ed25519, ротация ключей \\
цепочка, дерево, снапшот, защита от отката & реестр аудита Curator & лог прозрачности, Merkle, снапшоты, высшая принятая версия \\
план сессии, digest, только сужение & curator-agent-runtime & скомпилированный план, эпоха плана, делегирование вниз \\
receipt, грант как потолок & curator-agent-runtime & execution-receipt, call plan, host assurance \\
аренда, fencing, чекпоинт & ax & lease record, эпоха аренды, checkpoint record \\
передача сессии, конвертация & ax & handoff, force takeover, clone между агентами \\
задача, оценка, класс, ревью & борд задач & README и progress, конфигурация спавна \\
контракт задачи, версии & борд задач, реестр & README с критериями как тело контракта; версии как записи журнала \\
ключ идентичности, ротация & реестр аудита Curator & набор пинов, ротация через перекрытие, отзыв удалением \\
скилл remote-shell, пины & Curator & closure, артефакты с хешами, провайдеры \\
\bottomrule
\end{tabular}
\end{center}

\section{Curator: скиллы, closure и реестр}

Curator есть протокол пакетов для агентских скиллов: единица распространения содержит инструкции для модели, провайдеров инструментов и артефакты, всё с идентичностью по содержимому. Компилятор разрешает closure: полный набор пинов, из которого собирается план сессии. Два свойства важны для протокола.

Первое: runtime не поставляет инструментов, инструменты приходят как скиллы. Скилл удалённого shell есть обычный пакет с провайдерами (локальный, по SSH, в контейнере) и артефактами (бинарь демона под каждую платформу с хешами). Замена провайдера или демона это изменение хеша в closure, видимое в диффе плана. Это ровно та форма, в которой публичный клиент получает две свои привязки.

Второе: реестр аудита. Записи с подписями Ed25519 поверх канонического JSON, идентификатор ключа как префикс хеша ключа, ротация ключей через перекрытие, лог прозрачности с непрерывными номерами и хешем предыдущей записи, дерево Меркла, подписанные снапшоты с корнем, длиной и головой, защита от отката через высшую принятую версию как долговременное состояние, федерация реестров с правилом deny-wins, офлайн-бандлы с контрподписью при импорте. Это готовый журнал уровня L0. Ему не хватает четырёх вещей, и все они добавляются, а не переписываются: типы записей для задач, сертификат кворума как порог эндорсментов (формат уже хранит эндорсменты с их конвертами подписей), доказательства включения для светлых клиентов, секвенсер как роль сервиса и сертификат остановки как запись с особым правилом принятия.

Инвентарь ограничений Curator честно помечает на macOS ограничение файловой системы и сети как недоступное в непривилегированном портативном режиме. Демон среды поэтому есть verified host execution provider в терминах Curator: отдельно установленный доверенный компонент с точным хешем бинаря, permit и receipt. Seatbelt живёт в нём, а не в портативном режиме.

\section{Runtime агента}

Runtime разделён на три части: компилятор, который строит план сессии заранее и офлайн; брокер, который владеет консолью и границей исполнения; ядро, которое обрабатывает вывод модели и не имеет ни ключей, ни файловой системы. Четыре свойства runtime протокол берёт как есть.

План сессии скомпилирован, неизменяем в течение своей эпохи (это первое из трёх значений слова эпоха) и имеет digest; каждый запрос, вызов и результат несёт идентификатор эпохи и digest плана, и результаты из прошлой эпохи не принимаются. Протокол связывает digest плана с записью назначения, и среда сверяет его на каждом вызове.

Делегирование только сужает: план дочерней сессии есть подмножество родительского, ни одного инструмента, гранта или лиза сверх. Протокол поднимает это правило на уровень задач.

Грант есть потолок, план вызова есть применённый профиль, receipt есть наблюдённое. Для shell точность всегда неизвестна, поэтому shell работает на потолке; отсюда требование, чтобы потолок был санбоксом задачи, а не машиной. Схема receipt уже содержит digest плана, digest плана вызова, идентичность провайдера, отчёт об ограничении, дайджесты входа и выхода.

Доверенный путь: компонент, обрабатывающий вывод модели, не отвечает на одобрения. В протоколе тот же принцип: оркестратор, будучи моделью, не подписывает кворум своих записей.

Host assurance сообщается как полное, частичное или отсутствующее, и вызывающий, требующий полного, отказывается. Это инвариант I9 в готовом виде.

\section{ax: сессии, аренды, передача}

ax есть менеджер сессий агентов: у логической сессии ровно один активный владелец, реплики только читают, владение оформлено арендой с эпохой (второе значение слова) и случайным fencing-токеном, события владельца несут эпоху, пиры отвергают события младшей эпохи, но сохраняют их в ветке расхождений. Плановая передача требует доказательства покоя: ввод заблокирован, foreground и background idle, ноль открытых дочерних процессов и хендлов, стабильный дайджест хранилища. Принудительный захват есть отдельное явное действие. Чекпоинты адресуются по содержимому и служат валидированной базой передачи.

Клон между агентами (например, из одной агентской среды в другую) создаёт новую логическую сессию и не переносит аренду, владение бордом, одобрения, учётные данные и полномочия незавершённых операций. Это утилита конвертации сессий, о которой говорит протокол, и она уже воплощает правило: с сессией не путешествует авторитет. Общий принцип ax, транскрипты и записи есть content authority, а индекс есть перестраиваемая проекция, есть то же самое, что журнал истина, сессия кэш.

Удалённая точка входа ax это RPC по стандартным потокам поверх SSH без открытых портов; демон среды следует тому же образцу.

\section{Борд задач}

Борд живёт в git-репозитории проекта как дерево EPIC, STORY, TASK. У задачи есть README с описанием, областью, пришпиленной к коммиту, и критериями приёмки, и progress со статусом, флагом обязательности ревью, классом, оценкой в единицах сложности, блокировками и чеклистом; ресурсы задачи хранят логи запусков, вердикты ревью и патчи. Конфигурация спавна задаёт максимальную параллельность, изоляцию в копиях рабочего дерева с командами проверки, базовую ветку интеграции и потолки моделей по ролям разработчика и ревьювера.

Это почти содержимое журнала минус подписи, кворум и репликация. Оценка задачи есть цена, флаг ревью есть класс, потолки моделей по ролям есть принимаемые модели манифеста, команды проверки есть проверки из профиля контракта. Не хватает одного: версионирования критериев. README задачи сегодня правится как обычный файл; в журнале его изменение должно становиться новой версией контракта с отдельной авторизацией. Материализация журнала как борда в git, описанная в главе 5, есть просто признание того, что борд уже есть.

\section{Что берётся как есть и что расширяется}

Как есть: подписи, каноническая форма, цепочка, снапшоты, защита от отката; план сессии и его digest; делегирование только вниз; receipts; аренды с fencing; чекпоинты; клон без переноса полномочий; формат борда.

Расширяется: реестр до журнала с типами записей, кворумом, доказательствами включения и секвенсером; спавн борда до оркестратора с планировщиком по ёмкости и областям; скилл shell до восстановленной сессии с подписью на машине владельца.

Новое: демон среды, секвенсер журнала, протокол ноды. Следующая глава описывает их.

\chapter{Что нужно построить}

\begin{warnbox}[Глава заменена роадмапом экосистемы]
Состояние на 29 сентября 2026 года. Что строится и в каком порядке, теперь задают роадмап экосистемы (вики, \texttt{roadmap/ecosystem-roadmap.ru.md}, черновик v4: этапы M0--M8, треки PC, KP и CM) и спецификации модулей: конфигурация процесса (curator-playbook, трек PC), хост сессий (agent-session-host, этап SH), главный цикл узла (curator-keeper, трек KP), доверие (curator-trust), подписанные сообщения и аренды (waggle, трек CM). Распределённый пул идёт треком AR в curator-agent-runtime после этапов M4 и M7. Демон среды, секвенсер журнала и вход в санбокс пока описаны только в работе. Минимальный runtime отложен вместе с собственным агентским циклом (curator-keeper, §7.7). Текст главы сохранён как история; где он расходится с роадмапом или спецификацией модуля, правы они.
\end{warnbox}

Три новых компонента и три расширения существующих. Для каждого: ответственность, интерфейс в набросках и то, какие инварианты он обеспечивает. Наброски интерфейсов здесь не нормативны; нормативные документы со схемами и векторами соответствия пишутся уровнем ниже.

\section{Демон среды}

Демон среды есть провайдер исполнения на хосте: единственный процесс на среде, который принимает команды извне. Он принадлежит проекту, его хеш опубликован в манифесте, он не слушает портов и работает по стандартным потокам поверх SSH.

Ответственность: проверка ключа ноды и записи назначения; CA среды и выпуск сертификатов на задачу с принудительной командой; приём соединений (OpenSSH с сертификатами на L0, встроенный SSH-сервер в самом демоне для публичного роя); создание и уничтожение санбоксов; исполнение команд в санбоксе с лимитами и receipts; восстановленные сессии; фоновые процессы и spill; построение интеграционного кандидата; атомарное продвижение головы; операции раннера; светлый клиент журнала с защитой от отката; безопасный режим по сертификату остановки; инвентарь ограничений.

\begin{exbox}[{Поверхность демона, набросок}]
\begin{lstlisting}
permit.verify(record, proof)       -> ok | reject(reason)
cert.issue(assignment)             -> ssh_cert(principal=task, ttl=window, force_command)
sandbox.open(task, plan_digest)    -> sandbox_id
sandbox.close(sandbox_id)
exec(sandbox_id, session, cmd, limits) -> receipt
bg.start / bg.status / bg.log / bg.kill
files.read_range / files.apply_patch
candidate.build(task, R)           -> {M, I}
main.cas(M, I, integration_record) -> ok | stale(M_now)
runner.call(op_id, idem_key, op, args, record) -> outcome(done|rejected|unknown)
runner.reconcile(op_id)            -> outcome
stop.apply(certificate)            -> safe_mode
inventory()                        -> {fs: full, net: full, mem: partial, ...}
\end{lstlisting}
\end{exbox}

Обеспечивает инварианты I1 (включая проверку актуальности на границе исполнения), I3, I6 в части среды, I8 в части высшей версии, I9, I10.

Платформенные бэкенды: macOS через Seatbelt и клоны APFS с группой процессов и сторожем; Linux через bubblewrap с cgroups, с возможностью gVisor или Firecracker для чужих агентов. Бэкенд выбирается конфигурацией среды и отражается в инвентаре.

Планировщик слотов: очередь назначений, слоты как ёмкость, приоритет по эпохе и оценке. Egress-прокси: локальный HTTP-прокси с allowlist из манифеста, единственный сетевой выход из санбокса.

\begin{warnbox}[{Четыре места, где легко ошибиться}]
Первое: дать агенту любой путь в shell хоста, хоть на секунду, хоть для отладки. Второе: построить кандидата, а одобрение или проверки взять от исходной ревизии. Третье: принять снапшот журнала без проверки высшей версии. Четвёртое: повторить операцию раннера с исходом не установлен без сверки. Каждая ошибка обнуляет соответствующий инвариант, и первые три из них исполняемая модель (раздел 11.8) находит как нарушения.
\end{warnbox}

\section{Секвенсер журнала}

Секвенсер расширяет реестр аудита до журнала. Ответственность: приём подписанных записей от оркестратора, воркеров и кор-нод; присвоение порядка; проверка правил валидаторов на уровне сервиса (для ускорения, но не вместо подписей); сбор эндорсментов кор-нод до порога кворума; выпуск снапшотов; доказательства включения; сертификаты доступности; сертификат остановки как запись с особым правилом.

\begin{exbox}[{Расширение реестра, что добавляется}]
\begin{itemize}
\item схемы записей по типам из главы 5 вместо записи аудита артефакта, включая версии контракта, зависимости, кандидата и адресованные к нему одобрения и аттестации с тремя исходами;
\item поле кворума: список эндорсментов кор-ключей, запись совершена при достижении порога уровня;
\item метод \texttt{inclusion\_proof(entry)} и снапшот с корнем для светлых клиентов;
\item метод \texttt{availability(hash)}: реплики подтверждают хранение, сертификат прикладывается к записи;
\item тип записи остановки, принимаемый с одной подписью кор-ключа, и возобновления, принимаемого с кворумом;
\item на L2: интерфейс к движку консенсуса вместо собственного порядка.
\end{itemize}
\end{exbox}

Обеспечивает I4 (правила исполняются как детерминированные предикаты), I7 (эпоха и токен в записях), I8 (остановка и версии), и хранит материал для I2, I3, I5.

\section{Протокол ноды}

Протокол ноды есть единственный способ для контрибьютора говорить с проектом, одинаковый для обеих форм клиента.

\begin{exbox}[Сообщения]
\begin{lstlisting}
join(manifest_url, node_key, axes, identity_proofs) -> plan, authority_diff
heartbeat(capacity_left, reset_window, cost_class, slots_wanted)
assign(task, role, plan_digest, sandbox_endpoint)   <- from project
session.open / exec / close                 -> via environment daemon
submit(task, R, B, receipts_head, signature)
review.request(task, I_or_R)                <- from project
review.submit(task, verdict, findings_hash, signature)
attest.request(task, I) / attest.submit(task, I, result, signature)
leave()
\end{lstlisting}
\end{exbox}

Присоединение компилирует план из манифеста проекта и показывает разницу полномочий. Heartbeat несёт остаток ёмкости, окно сброса и класс стоимости. Подписи ставятся брокером на машине ноды.

Обеспечивает I5 в части объявления осей, I6, и даёт секвенсеру данные для I2.

\section{Резолверы идентичности}

Слой идентичности (глава 8) в MVP не входит, но интерфейс стоит зафиксировать сейчас, чтобы протокол ноды и типы записей его не исключали. Резолвер есть плагин с одним методом: по строке идентичности вернуть набор отпечатков ключей и свидетельство. Первые три: TXT-запись домена под DNSSEC, well-known по HTTPS, ключи аккаунта на хостинге кода.

\begin{exbox}[Интерфейс резолвера]
\begin{lstlisting}
resolve(identity) -> {fingerprints: [..], evidence_hash, fetched_at, resolver_id}
                   | unavailable(reason)
\end{lstlisting}
Запись подтверждения идентичности содержит идентичность, класс доказательства, хеш свидетельства и окно; она совершена, когда $t$ валидаторов, каждый резолвивший сам, записали одинаковый хеш. Запись сертификации ключа содержит ключ ноды, подписанный ключом идентичности, с началом и концом действия. Файл допустимых подписантов git для среды выводится из этих двух типов записей.
\end{exbox}

\section{Скилл удалённого shell}

Скилл есть пакет Curator: инструкции для модели (минимальные: у вас есть shell в рабочем окружении), провайдер, который говорит с демоном среды, и артефакты. Провайдер реализует восстановленную сессию (снимок состояния по отдельному дескриптору, восстановление перед командой), именованные сессии с сериализацией мутирующих вызовов, фоновые процессы по дескриптору, spill больших выводов, отмену по группе процессов, переподключение без потери состояния сессии.

Отличие от демо на одноразовой машине: демон принадлежит проекту, не заливается контрибьютором, а пинится по хешу из манифеста. Провайдер проверяет, что подключился к тому демону, идентичность которого объявлена, и это доверие к проекту, а не верификация.

Вторая привязка публичного клиента, подпись ревизий, есть операция брокера: подписать хеш, который совпадает с хешем в журнале задачи.

\section{Минимальный runtime}

Публичный клиент есть runtime с планом из двух привязок и командой присоединения. Установка одной командой: ставится runtime, по адресу манифеста компилируется план, закрывается closure с пинами скилла shell, скилла подписи и хеша демона, показывается разница полномочий, и это вся установка. Работа над тем, чтобы весь набор ставился одной командой и был бесшовным для контрибьютора, уже идёт.

\section{Оркестратор поверх спавна борда}

Оркестратор есть кор-нода в роли: агент, который читает цели, пишет задачи в борд, отправляет записи в секвенсер и назначает ноды. Спавн борда уже умеет запускать роли разработчика и ревьювера в копиях дерева с командами проверки; оркестратор заменяет локальный запуск назначением на ноду через протокол, добавляет планировщик по ёмкости и областям, заморозку кандидатов и случайный выбор ревьювера, меню при превышении $N$, аренду с heartbeat и чекпоинты.

\section{Два эксперимента с первого дня}

Соблюдение протокола и полезность роя есть разные вопросы, и они проверяются по-разному.

\term{Соблюдение протокола.} Реализация подвергается сбоям и нарушителям, воспроизводимым программно, без участия моделей: рабочая нода, выходящая за write scope; запись оркестратора с устаревшей арендой; подмена кандидата после аттестации; изменение критериев после ревью; повторная доставка записи об интеграции; операция раннера с потерянным исходом; смерть оркестратора с удалением его сессии. Критерий: ни одно нарушение не приводит к эффекту на среде или продвижению головы, а восстановление проходит по журналу. Последний сценарий, остановить оркестратора, удалить его сессию и проверить, продолжит ли преемник из журнала и доступных артефактов, есть прямая проверка инварианта I7 и утверждения о достаточности внешнего состояния.

\term{Полезность роя.} Единица сравнения есть заранее определённая внешняя задача или завершённая функция проекта с одинаковыми критериями приёмки, а не внутренние задачи оркестратора: он сам декомпозирует работу и может менять их число, не меняя объёма результата. Три условия при сравнимых ресурсных ограничениях: одиночный агент; параллельный процесс без протокола, агенты одного владельца с общим репозиторием и блокировками; развёртывание L0. Стоимость включает inference, исполнение проверок, занятость среды и человеческое вмешательство. Диагностические метрики собираются из журнала и receipts без дополнительной инструментовки: принятые задачи на миллион токенов, время от авторизации до интеграции, ожидание слотов, блокировки областей, повторная работа после смены головы, восстановление после потери агента, доля пересмотренных отклонений, доля откатов через $k$ эпох, часы человеческого ревью на принятую задачу.

Дашборд поверх журнала есть отдельная маленькая задача; важно, чтобы данные для него существовали с первого совершённого назначения. Ориентир масштаба для первого числа: эксперимент, в котором три личные подписки за три дня совместно формализовали крупный результат.

\section{Исполняемая модель фрагмента}

Инварианты и правила валидаторов детерминированы, и ключевой фрагмент протокола допускает исполняемую модель. К третьей редакции работы приложена модель фрагмента кандидат, аттестации, запись об интеграции, атомарное продвижение головы, проверенная исчерпывающим перебором состояний в заданных границах.

\begin{exbox}[Что в модели]
Две задачи: $A$ с полным поведением (контракт может менять версию, записи могут нести устаревшую эпоху) и $B$, которая только интегрируется и тем сдвигает голову. Состояние: голова, эпоха аренды и для каждой задачи наличие ревизии, база кандидата, аттестация (адрес и версия контракта), версия контракта, признак интеграции, запись об интеграции с числом доставок. Действия: представить ревизию; построить кандидата над текущей головой; аттестовать кандидата; изменить версию контракта; сменить эпоху; предложить запись с текущей или устаревшей эпохой; доставить запись, в том числе повторно. Границы: не больше двух версий контракта, двух эпох, двух доставок на запись. Правила при доставке: голова равна базе записи; запись несёт текущую эпоху; аттестация ссылается на текущую версию контракта; аттестация адресована ровно к кандидату. Свойство: голова продвигается к кандидату, построенному над головой в этот момент, с аттестацией к этому кандидату и текущему контракту, по записи текущей эпохи, и ни одна запись не применяется дважды.
\end{exbox}

При всех четырёх правилах достижимо 1090 состояний, свойство выполнено во всех. При отключении любого одного правила перебор находит контрпример: без проверки головы одна запись применяется дважды при повторной доставке, а также голова продвигается к кандидату, построенному не над текущей головой, после интеграции другой задачи; без проверки эпохи запись устаревшей аренды продвигает голову; без проверки версии контракта аттестация под прежними критериями принимается после их изменения; без привязки к кандидату аттестация, сделанная на базе до сдвига головы, переносится на кандидата над новой головой. Последний контрпример есть ровно сценарий переноса одобрения по patch-id, и именно он вывел это правило из гарантии.

\begin{warnbox}[Границы результата]
Это проверка ограниченного пространства сценариев, а не доказательство для всех допустимых исполнений: границы заданы явно, среда и валидаторы в модели честны, суждение и содержимое кода не моделируются. Что модель даёт реализации: её трассы становятся векторами соответствия демона среды и секвенсера, и первый набор векторов получается запуском модели, а не ручным написанием. Следующий шаг есть та же модель в языке спецификаций с механической проверкой и распространение на аренду и остановку.
\end{warnbox}

\chapter{План на двенадцать недель}

\begin{warnbox}[Глава заменена роадмапом экосистемы]
Состояние на 29 сентября 2026 года. План на двенадцать недель не исполняется. Порядок работ задаёт список приоритетов роадмапа экосистемы (вики, \texttt{roadmap/ecosystem-roadmap.ru.md}, черновик v4, §2b): первым идёт модуль хоста сессий и на нём waggle CM1, вторым конфигурация процесса PC0--PC1 с первыми демонстрационными плейбуками, дальше этапы M0--M8 и треки CM, PC, KP и AR. Сроки и готовность каждого шага ведут треки в репозиториях модулей. Текст главы сохранён как история; где он расходится с роадмапом, прав роадмап.
\end{warnbox}

Цель: через три месяца работающий MVP, который соединяет существующую инфраструктуру в продукт и готов к обкатке на реальных проектах. Пилот не есть отдельный этап: метрики встроены в MVP, и обкатка начинается в день готовности. Уровень доверия первого развёртывания L0, с определением кворума как параметром, чтобы L1 и L2 не требовали переписывания.

\section{Три блока}

\begin{center}
\begin{ganttchart}[
  hgrid, vgrid={*{3}{draw=none},dotted},
  x unit=2.6mm, y unit chart=6.5mm, y unit title=6.5mm,
  title label font=\footnotesize\sffamily, bar label font=\footnotesize,
  bar/.style={fill=paleblue,draw=accent}, bar height=0.55,
  group/.style={fill=accent,draw=accent}, group label font=\footnotesize\sffamily\bfseries,
  milestone/.style={fill=accentamber,draw=accentamber}, milestone label font=\footnotesize
]{1}{36}
\gantttitle{Недели 1--4}{12}\gantttitle{Недели 5--8}{12}\gantttitle{Недели 9--12}{12}\\
\ganttgroup{Блок 1: исполнение}{1}{12}\\
\ganttbar{Демон среды, macOS и Linux}{1}{10}\\
\ganttbar{Скилл shell, восстановленная сессия}{3}{10}\\
\ganttbar{Минимальный runtime, подпись}{5}{11}\\
\ganttmilestone{Агент проходит полный цикл}{12}\\
\ganttgroup{Блок 2: журнал}{13}{24}\\
\ganttbar{Секвенсер поверх реестра}{13}{21}\\
\ganttbar{Светлый клиент, откат, включение}{16}{22}\\
\ganttbar{Кандидат, patch-id, CAS}{17}{23}\\
\ganttmilestone{Интеграция только по QC}{24}\\
\ganttgroup{Блок 3: рой}{25}{36}\\
\ganttbar{Аренда, планировщик, области}{25}{33}\\
\ganttbar{Случайный ревьювер, три оси}{27}{33}\\
\ganttbar{Остановка, join, метрики}{29}{35}\\
\ganttmilestone{MVP, обкатка на проектах}{36}
\end{ganttchart}
\end{center}

\subsection{Блок 1, недели 1--4: исполнение}

Демон среды как провайдер исполнения на хосте: Seatbelt и bubblewrap, санбокс на задачу с клоном дерева, защищённые пути только на чтение, receipts на команду, инвентарь, планировщик слотов, egress-прокси, CA среды с сертификатами на задачу и принудительной командой поверх OpenSSH (встроенный SSH-сервер в демоне откладывается до публичного роя). Скилл shell с восстановленной сессией, именованными сессиями, фоном, spill и переподключением. Минимальный runtime с планом из двух привязок и подписью на машине владельца.

Критерий приёмки: один агент проходит полный цикл на реальном репозитории проекта через среду, а не локально: клонирует, правит, собирает, тестирует, представляет подписанный коммит; транскрипт структурно неотличим от локального прогона; ни одной команды вне санбокса; receipts на каждую команду.

\subsection{Блок 2, недели 5--8: журнал}

Секвенсер поверх реестра аудита: типы записей, включая версии контракта и зависимости, кворум как порог эндорсментов, снапшоты, доказательства включения, сертификаты доступности, сертификат остановки. Среда как светлый клиент: открытие санбокса только по назначению с QC, защита от отката, проверка актуальности на границе исполнения. Интеграционный кандидат до финального одобрения, атомарное продвижение головы, операции раннера с ключами идемпотентности и тремя исходами. Борд как содержимое журнала: переходы как коммиты борда с QC.

Критерий приёмки: ни один санбокс не открывается и ни одна интеграция не проходит без записи с QC; одобрение и аттестации, адресованные не к интегрируемому кандидату или к прежней версии контракта, отвергаются; попытка интегрировать при ушедшей вперёд голове приводит к пересборке кандидата с повторным ревью; повторная подача старого снапшота отвергается; операция с исходом не установлен не повторяется без сверки. Все трассы исполняемой модели проходят как векторы.

\subsection{Блок 3, недели 9--12: рой}

Аренда оркестратора с fencing и heartbeat поверх спавна борда; планировщик по ёмкости, окнам квот, слотам и областям; заморозка кандидатов, случайный ревьювер, три оси; меню при превышении $N$; сертификат остановки и возобновление; команда присоединения по манифесту; дашборд метрик.

Критерий приёмки: три и больше контрибьюторов с ёмкостью разного происхождения одновременно ведут задачи одной цели; ревьювер никогда не совпадает с продьюсером по трём осям; остановка одной подписью прекращает открытие санбоксов и интеграции; смена оркестратора по таймауту не теряет состояния; пять метрик показываются из журнала.

\section{Критический путь}

Демон среды и скилл shell. Всё остальное можно временно заменить заглушками: журнал секвенсером на одной ноде, оркестратора человеком, ревьювера вторым агентом на той же машине. Но без демона со санбоксом и без shell, неотличимого от локального, нет ничего, что можно показать. Поэтому первые две недели уходят только туда, и туда же направляется большинство агентов.

\section{Что отложено без потери формы}

Уровни L1 и L2: кворум есть параметр, и реализация L0 не закрывает путь. Переносимая репутация между проектами: аттестации пишутся с первого дня, взвешивание позже. Экономика и учёт вклада: данные пишутся (оценка, исход, вестинг), правила раздела пулов считаются потом. Полировка передачи сессий между агентами: плановая передача через ax уже есть, тонкая настройка позже. Верификатор как отдельная машина: сначала одна среда, вторая при первой возможности. Реестр федерации: сначала манифест по адресу, реестр когда проектов больше одного. Слой идентичности: типы записей и поле доказательств в join закладываются сейчас, резолверы и подтверждение кворумом валидаторов в первом релизе после MVP.

\section{Как двигаться с агентами}

Скорость достигается тем же циклом, которым проект уже живёт: спецификация впереди реализации, с нормативными формулировками и векторами соответствия, затем агенты реализуют против векторов параллельно. Порядок:

\begin{enumerate}
\item Неделя 1: нормативные документы для демона и для типов записей журнала. Первый набор векторов соответствия получается запуском исполняемой модели фрагмента интеграции, а не ручным написанием; вручную дописываются векторы для санбокса и внешних эффектов. Это единственная неделя, где узкое место люди.
\item Со второй недели: агенты реализуют бэкенды, скилл, секвенсер и клиент против векторов, каждая задача через борд с ревью.
\item Люди держат два места, где нужна голова: модель полномочий (что считать областью, как компилировать план из области) и семантику интеграционного кандидата (что делать при неравенстве patch-id, как строить очередь).
\end{enumerate}

Когда демон и скилл готовы, рой начинает строить сам себя: следующие задачи проекта идут через протокол, и первые метрики появляются не на чужом проекте, а на своём.

\section{Риски}

\begin{center}
\small
\begin{tabular}{p{4.6cm} p{8.8cm}}
\toprule
Риск & Что делаем \\
\midrule
условия вендоров по подпискам в стороннем runtime & ёмкость любого происхождения: ключи API и локальные модели как равноправные источники; подписки там, где допустимо; проверка условий до публичного запуска \\
ёмкость одной среды & слоты и очередь с первого дня; вторая машина как реплика и верификатор при первой возможности \\
ошибки в санбоксе на macOS & инвентарь честно объявляет частичное; никаких общих записываемых каталогов; всё ценное за операциями раннера; тест на побег как часть векторов \\
стимулы ревью (отклонять на всякий случай) & измеряем долю пересмотренных отклонений; правило оплаты раундов и споров настраивается по данным \\
шторм перебазирований и повторных ревью & дизъюнктные области как правило валидатора; декомпозиция по областям; стоимость повторного ревью кандидата измеряется, перенос одобрения возможен только как расширение с условием по прочитанным путям \\
терминология (три значения слова эпоха) & глоссарий в спецификациях; в коде разные имена: plan epoch, lease epoch, orchestrator epoch \\
\bottomrule
\end{tabular}
\end{center}

\chapter{Открытые вопросы}

Решённое отделено от нерешённого честно. Ниже вопросы, по которым механика ясна, а политика или детали нет, и вопросы, по которым нет и механики. Для каждого текущее состояние мысли, чтобы новый участник не начинал с нуля. С версии 0.4 там, где вопрос перешёл в спецификацию модуля, в конце раздела сказано, где он живёт сейчас.

\section{Пересечение областей}

Несколько оркестраторов с дизъюнктными областями работают. Задача, затрагивающая две области, принадлежит аренде, покрывающей объединение, либо создаётся кор-нодой напрямую. Не решено: кто и по какому правилу решает, что область нужно объединить, и как часто это будет происходить на реальном монорепозитории. Ожидание: редко, если декомпозиция идёт по модулям с явными интерфейсами; проверяется на данных.

\term{Где вопрос сейчас.} Владение областями задают аренды областей waggle (\texttt{spec/waggle.md} v5.2, §12): ядро проверяет аренду при запуске работы, при интеграции и при записи борда, а область переходит к другому оркестратору через \texttt{coord.handoff-offer} с большей эпохой аренды. Для нескольких операторов на одном борде решён один принимающий писатель с fencing по эпохе (трек CM2); как он выбирается, по владению процессом или выбором участников, открыто (вики, \texttt{open-questions.md}).

\section{Стимулы ревью}

Вестинг для одобрившего ревьювера толкает к отклонению на всякий случай, и пропускная способность падает. Правило неоплаты проигранных в споре отклонений давит в обратную сторону. Это классическая задача peer prediction, у неё нет закрытого решения, и в протокол оно не заложено. Измеряется долей пересмотренных отклонений; правила оплаты раундов настраиваются по данным.

\term{Где вопрос сейчас.} Измерять пропуски ревьювера предлагает калибровка посевом дефектов с охранным условием \texttt{review.recall} (curator-playbook, \texttt{spec/process-configuration.md} v4.8, §6.16). Лишние отклонения она не измеряет, и вопрос о стимулах остаётся открытым.

\section{Семантически вредная задача внутри своей области}

Ограничена областью, бюджетом, независимым ревью, вестингом и остановкой, но не устранена. Возможные усиления: детерминированное правило для дифов деструктивной формы (удаление больше заданной доли строк, удаление файлов вне объявленной области) с требованием дополнительного одобрения кор-ноды; проверка семантического соответствия контракта задачи родительской цели вторым ревью для обычного класса. Версионирование контракта закрывает тихое ослабление критериев, но не вредную задачу с честными критериями. Оба усиления удорожают цикл, и вопрос о пороге открыт.

\section{Перенос одобрения между кандидатами}

Базовая гарантия требует повторного ревью после каждой пересборки кандидата. Расширение вне гарантии допускает перенос при условии, что дельта базы не пересекается ни с write scope задачи, ни с множеством путей, прочитанных продьюсером по receipts. Не решено: достаточно ли множества прочитанных путей как приближения контекста ревью, и как часто условие будет выполняться на реальных очередях. Проверяется на данных пилота: доля пересборок, при которых условие выполнено, и доля дефектов, найденных повторным ревью там, где оно было бы пропущено.

\section{Механическая проверка}

Исполняемая модель фрагмента интеграции проверена перебором в заданных границах. Не сделано: та же модель в языке спецификаций с механической проверкой, распространение на аренду, остановку и внешние эффекты, и доказательство свойств для всех допустимых исполнений, а не для ограниченного пространства сценариев.

\term{Где вопрос сейчас.} Исследование формальных методов записано в роадмап экосистемы (вики, \texttt{roadmap/ecosystem-roadmap.ru.md}, §5: проверка моделей для критичных протоколов, например аренд waggle), а блок проверяемых свойств процесса с экспортом в TLA+ или Quint решено делать после PC1 (curator-playbook v4.8, §17.3, п.~17). Модель фрагмента интеграции в эти планы пока не включена.

\section{Аттестация среды}

Среда доверена в отношении целостности. На Linux доступны аппаратно аттестованные анклавы и microVM, на macOS нет ни того, ни другого для произвольной нагрузки. Верификатор на второй машине переводит тесты из подписи под ответом той же машины в независимое исполнение, но не закрывает подмену дифа, показанного ревьюверу. Возможный ход: ревьювер получает диф не из санбокса среды, а по хешу коммита из собственной копии объектов git, что сводит доверие к целостности хешей. Не спроектировано.

\section{Экономика}

Учёт вклада описан правилами, но ставки (цена единицы сложности, множители для ёмкости и для человеческого времени, доля ретроактивного пула) не назначены, а форма конвертации (репутация, деньги, доля) зависит от проекта. Для доверенной группы задача решаема как контракт плюс время плюс ограниченная дискреция. Для открытого анонимного роя пропорциональные доли не собираются в принципе, и там остаются репутация и баунти. Экономика вынесена в отдельную работу.

\section{Пул сред}

Одна авторитетная среда и реплики по снимкам описаны, но протокол синхронизации снимков между репликами и правило, какая реплика обслуживает какой санбокс, не специфицированы. Ожидание: реплики получают снимок дерева на момент назначения по хешу коммита из общего хранилища объектов, и любая реплика может обслуживать любой санбокс; кандидат строится и голова продвигается только на авторитетной.

\section{Федерация и реестр}

Манифест проекта как источник плана контрибьютора описан. Реестр манифестов, взвешивание чужих аттестаций репутацией выдавшего проекта и правила деления ёмкости между проектами описаны на уровне принципа. Всё это нужно, когда проектов больше одного, и не раньше.

\term{Где вопрос сейчас.} Общение с другими организациями есть этап waggle CM4: CA партнёров, инспекторы сообщений по умолчанию, федерация XMPP, шлюз A2A, квоты (\texttt{spec/waggle.md} v5.2, §17). Допуск внешних участников есть стадия L2 в curator-trust (\texttt{spec/trust.md} v0.4, §13): репутация удешевляет проверку, но доверия не даёт. Реестр манифестов и деление ёмкости между проектами ни одна спецификация пока не описывает.

\section{Идентичность: вес псевдонимов и хрупкость провайдеров}

Механика доказательств идентичности ясна, политика нет. Не решено: какой вес и какую глубину ревью давать псевдонимным нодам, чтобы не отрезать предложение ёмкости и не открыть дешёвые множественные личности; как часто переподтверждать доказательства, когда провайдер закрыл доступ к API или удалил публикацию; считать ли ключи аккаунта на хостинге кода доказательством того же класса, что домен под DNSSEC. Ожидание: веса настраиваются по данным первых проектов, а хрупкость провайдеров гасится требованием двух доказательств разных классов для веса выше псевдонимного.

\term{Где вопрос сейчас.} Классы ключей и уровни их гарантии (обычный, привязанный к устройству, аттестованный) задаёт waggle (\texttt{spec/waggle.md} v5.2, §4.5). Кто считается человеком, какие ключи требуют его присутствия и как ключи одного человека ручаются друг за друга, задаёт curator-trust (\texttt{spec/trust.md} v0.4, §3). Доказательства идентичности через резолверы и их веса в эти спецификации пока не перенесены.

\section{Что не решено намеренно}

Полная децентрализация без единой среды: техническая возможна через аттестованное или повторное исполнение, но не нужна, пока доверие к среде решается социально, выбором проекта. Блокчейн как расчётный слой: имеет смысл только когда проекты начнут платить контрибьюторам деньгами или потребуется стейк как барьер вместо цены ёмкости. Оба вопроса отложены с формулировкой, а не забыты.

\appendix
\chapter{Глоссарий}

\begin{description}[leftmargin=0pt,labelindent=0pt,itemsep=3pt,font=\sffamily\bfseries]
\item[Аренда] Запись, закрепляющая роль оркестратора за кор-нодой на срок, продлеваемый heartbeat; несёт эпоху, fencing-токен, предшественника, причину, чекпоинт и область.
\item[Аттестация] Подписанный результат проверки одного свойства из профиля контракта для конкретного кандидата, версии контракта, политики и окружения; исход есть пройдена, провалена или не установлена.
\item[Адрес аттестации] Четвёрка: версия контракта, кандидат, версия политики проверки, digest окружения; аттестация со ссылкой на устаревшую часть недействительна.
\item[Верификатор] Реплика среды, независимо повторяющая критические прогоны тестов.
\item[Восстановленная сессия] Механизм shell, при котором каждый вызов есть отдельный процесс, а состояние сессии снимается и восстанавливается между вызовами.
\item[Доказательство идентичности] Двусторонняя привязка ключа идентичности к внешней сущности: заявление, подписанное ключом, и артефакт у провайдера с отпечатком ключа; подтверждается совпадением свидетельств у $t$ валидаторов.
\item[Журнал] Упорядоченная цепочка подписанных записей, единственное состояние авторитета.
\item[Запись] Единица журнала: хеш предыдущей, тип, тело или хеши тел, подписи.
\item[Инвентарь] Объявление среды о том, какие ограничения она обеспечивает полностью, частично или не обеспечивает.
\item[Интеграционный кандидат] Неизменяемое состояние $I$, построенное из ревизии $R$ над текущей головой $M$; предмет финального одобрения, проверок и интеграции.
\item[Исход проверки] Одно из трёх: пройдена, провалена, не установлена; сбой или таймаут инструмента дают не установлена.
\item[Контракт задачи] Версионируемое описание того, что задача должна получить и по каким правилам результат принимается: цель, write scope, зависимости, критерии, профиль свойств, оценка, класс, бюджеты. Изменение критериев создаёт новую версию с авторизацией.
\item[Ключ идентичности] Долгоживущий ключ владельца ноды, лежащий в доказательствах идентичности; сертифицирует короткоживущие ключи нод записью в журнале.
\item[Кор-нода] Узел с фиксированным ключом, ведущий журнал и проверяющий правила.
\item[Оркестратор] Прикладная роль планирования, удерживаемая кор-нодой по аренде.
\item[Патч-id] Хеш содержимого дифа, не зависящий от базы. Равенство patch-id означает то же изменение на другой базе, но не тот же смысл; поэтому перенос одобрения по patch-id не входит в гарантию протокола.
\item[Профиль свойств] Часть контракта: какое свойство требуется подтвердить, каким способом, для какого артефакта и при каких допущениях.
\item[План] Скомпилированный набор полномочий агентской сессии с digest.
\item[Принимающий писатель] Единственный участник, который принимает записи в борд процесса, когда его ведут несколько операторов; ограждён fencing по эпохе аренды. Направление решено 28 сентября 2026 года и идёт в трек waggle CM2.
\item[Рабочая нода] Агентский цикл контрибьютора с его ёмкостью и ключами; не доверена.
\item[Раннер] Доверенный компонент среды, исполняющий именованные операции с секретами и общими сервисами.
\item[Резолвер] Плагин, возвращающий по строке идентичности набор отпечатков ключей и свидетельство; DNS TXT, well-known, ключи аккаунта на хостинге кода и другие.
\item[Receipt] Подписанный факт исполнения одной команды: дайджесты входа, выхода и состояния, код выхода, ресурсы.
\item[Санбокс] Копия рабочего дерева и политика ограничения для одной задачи.
\item[Секвенсер] Сервис, упорядочивающий записи журнала.
\item[Сертификат кворума, QC] Набор подписей, достаточный для совершения записи на данном уровне доверия.
\item[Сертификат доступности] Подтверждение реплик о хранении объекта, прикладываемое к хешу в журнале.
\item[Сертификат остановки] Запись одной кор-подписи с эпохой и номером, переводящая среду в безопасный режим.
\item[Сертификат на задачу] Короткоживущий SSH-сертификат, который среда выпускает при совершении назначения: principal равен задаче, срок равен окну назначения, принудительная команда ведёт только в санбокс задачи.
\item[Слот] Право на одновременное исполнение; число слотов есть ёмкость среды.
\item[Среда исполнения] Авторитетная машина проекта с рабочим деревом, санбоксами, тестами и основной ветвью, возможно с репликами.
\item[Три оси] Административный домен, семейство модели, тулчейн; основа независимости участников.
\item[Уровни доверия L0, L1, L2] Одни инварианты, разное определение кворума: одна подпись; больше половины; $2f+1$ из $3f+1$.
\item[Харнес] Программа, в которой работает агент: Claude Code, Codex, Muse, харнес над локальной моделью. Собственного харнеса у системы нет, и независимость от харнеса есть принцип дизайна.
\item[Fencing-токен] Растущий номер аренды, по которому исполнители отвергают действия устаревшего владельца.
\item[Write scope] Множество путей, которые задача вправе изменять; удерживается с назначения до интеграции или отмены; предотвращает конкурирующую запись, но не семантические конфликты.
\item[Эпоха] В этой книге срок аренды оркестратора. В runtime агента: время жизни плана сессии. В менеджере сессий: эпоха аренды владельца сессии.
\end{description}

\section*{Карта имён}

Имена для людей операторы утвердили 25 сентября 2026 года (вики, \texttt{naming.ru.md}). Материалы для людей пользуются ими, а спецификации, код и эта книга пишут технические термины. Плейбук и waggle сохраняют свои имена; слово рой в этой книге описательное и именем для людей не служит.

\begin{center}
\small
\begin{tabular}{p{2.4cm} p{3.4cm} p{7.6cm}}
\toprule
Имя для людей & Технический термин & В этой книге \\
\midrule
пасека & вся система & Apiary: все узлы и процессы одного владельца; протокол книги есть её распределённый слой \\
точка & узел & одна машина; на ней работает кор-нода, рабочая нода или среда исполнения \\
Кипер & главный цикл & прямого аналога нет; главный цикл узла запускает корневого оркестратора процесса \\
улей & процесс & проект со своей средой, журналом и оркестратором \\
матка & корневой оркестратор & оркестратор (глава 7) \\
соты & борда & борд задач, в котором материализуется журнал (глава 5) \\
пчела & запуск & агентская сессия рабочей ноды по одному назначению \\
дикие пчёлы & внешние участники & рабочие ноды контрибьюторов (глава 8) \\
нуклеус & пакет процесса & прямого аналога нет; опубликованный плейбук, с которого начинается новый процесс \\
\bottomrule
\end{tabular}
\end{center}

\chapter{Инварианты}

\begin{enumerate}[label=I\arabic*.,leftmargin=2.6em,itemsep=3pt]
\item Ни одного эффекта на среде без совершённой записи авторизации, digest плана которой совпадает с digest плана вызова, и актуальность которой проверена на границе исполнения.
\item Полномочия только сужаются: план задачи есть подмножество capability, capability есть часть цели, ревизия не изменяет ничего вне write scope; условия приёмки задачи меняются только новой версией контракта с отдельной авторизацией.
\item Основная ветка продвигается только атомарной заменой головы $M$ на неизменяемого кандидата $I$, к которому привязаны финальное одобрение и аттестации, ссылающиеся на текущую версию контракта и политики.
\item Суждение не входит в консенсус: валидаторы проверяют только детерминированные предикаты над совершёнными записями и хешами содержимого.
\item Продьюсер, ревьювер и исполнитель проверок одной задачи независимы согласно объявленной политике; отсутствие подходящих участников не ослабляет политику автоматически; случайность выбора фиксируется после заморозки кандидатов и хеша кандидата.
\item Ключи контрибьютора не покидают его машину, и подпись выполняется на машине агентского цикла; секреты проекта существуют только в раннере и недоступны ни из одного санбокса.
\item Журнал есть единственное состояние авторитета, сессия модели есть кэш; записи с устаревшей арендой отвергаются и сохраняются как расхождение.
\item Любая кор-нода останавливает систему, кворум уровня возобновляет; среда хранит наивысшую принятую версию журнала и не принимает более раннюю.
\item Ограничение есть сообщаемый факт: среда объявляет, что она способна ограничить, и задача, требующая полного ограничения, отказывает на частичном.
\item Все обращения к среде через протокол и выданные ею результаты имеют проверяемое происхождение и связаны цепочкой receipts; протокол не удостоверяет полноту внутреннего контекста недоверенной рабочей ноды.
\end{enumerate}

\chapter{Чеклист приёмки MVP}

\section*{Исполнение}
\begin{itemize}
\item Ни одна команда агента не исполняется вне санбокса задачи; попытка выйти за политику отвергается ядром, а не инструментом.
\item Защищённые пути смонтированы только на чтение; изменить тест из обычной задачи невозможно.
\item Операции с ключами и устройствами недоступны из shell и проходят только через раннер по записи.
\item Вход только по сертификату на задачу с restrict и принудительной командой, от непривилегированного пользователя; транспортный ключ не совпадает с ключом подписи; после закрытия задачи сессии рвутся.
\item Каждая команда оставляет receipt; receipts задачи образуют цепочку; голова входит в ревизию.
\item Инвентарь ограничений опубликован и честен; задача с требованием полного ограничения отказывает на частичном.
\item Транскрипт агента через среду структурно неотличим от локального: cd и переменные переживают вызовы, фоновые процессы живут по дескриптору, большой вывод не убивает контекст, обрыв сети не теряет сессию.
\end{itemize}

\section*{Журнал}
\begin{itemize}
\item Санбокс открывается только по записи назначения с QC; интеграция проходит только по записи с QC.
\item Digest плана каждого вызова сверяется с записью назначения; несовпадение отвергается.
\item Старый снапшот не принимается; высшая принятая версия переживает перезапуск среды.
\item Хеш тела включается в журнал только с сертификатом доступности.
\item Кандидат строится над текущей головой до финального одобрения; одобрение и аттестации адресованы к кандидату, версии контракта, политике и окружению; при уходе головы кандидат перестраивается, ревью и проверки повторяются.
\item Изменение критериев приёмки создаёт новую версию контракта и требует авторизации; аттестации под прежней версией отвергаются.
\item Операция раннера имеет идентификатор и ключ идемпотентности; исход не установлен не повторяется без сверки; актуальность разрешения проверяется на границе исполнения.
\item Трассы исполняемой модели фрагмента интеграции проходят как векторы соответствия.
\item Все правила валидаторов проверяются без обращения к модели.
\end{itemize}

\section*{Рой}
\begin{itemize}
\item Ревьювер вычислим любым валидатором из заморозки, хеша ревизии и случайности; продьюсер, ревьювер и тестировщик различны по трём осям.
\item Задачи с пересекающимися write scope не исполняются одновременно.
\item Смена оркестратора по таймауту не теряет состояния; запись с устаревшим токеном отвергается и сохраняется как расхождение.
\item Одна подпись останавливает; возобновление требует кворума.
\item Присоединение контрибьютора занимает одну команду и показывает разницу полномочий.
\item Пять метрик показываются из журнала.
\end{itemize}

\chapter{Литература}

\begin{english}
\begin{enumerate}[label={[\arabic*]},leftmargin=2.6em,itemsep=2pt]
\item A. Grigorev, I. Oparin. Software Development by a Swarm of Autonomous Agents over a Single Execution Environment (in Russian). Relux Works, third edition, version 3.3, 2026.
\item L. Lamport, R. Shostak, M. Pease. The Byzantine Generals Problem. ACM TOPLAS 4(3), 1982.
\item M. Castro, B. Liskov. Practical Byzantine Fault Tolerance. OSDI, 1999.
\item D. Ongaro, J. Ousterhout. In Search of an Understandable Consensus Algorithm. USENIX ATC, 2014.
\item E. Buchman. Tendermint: Byzantine Fault Tolerance in the Age of Blockchains. M.Sc. thesis, University of Guelph, 2016.
\item M. Yin, D. Malkhi, M. K. Reiter, G. Golan Gueta, I. Abraham. HotStuff: BFT Consensus with Linearity and Responsiveness. PODC, 2019.
\item D. P. Anderson. BOINC: A System for Public-Resource Computing and Storage. Grid Computing Workshop, 2004.
\item A. Cheng, E. Friedman. Sybilproof Reputation Mechanisms. P2PEcon, 2005.
\item S. Micali, M. Rabin, S. Vadhan. Verifiable Random Functions. FOCS, 1999.
\item J. R. Douceur. The Sybil Attack. IPTPS, 2002.
\item R. C. Merkle. A Digital Signature Based on a Conventional Encryption Function. CRYPTO, 1987.
\item B. Laurie, A. Langley, E. Kasper. Certificate Transparency. RFC 6962, 2013.
\item Z. Newman, J. S. Meyers, S. Torres-Arias. Sigstore: Software Signing for Everybody. ACM CCS, 2022.
\item S. Torres-Arias et al. in-toto: Providing Farm-to-Table Guarantees for Bits and Bytes. USENIX Security, 2019.
\item G. Hoare. The Not Rocket Science Rule of Software Engineering. 2014.
\item Remote Execution API. github.com/bazelbuild/remote-apis.
\item S. Chen et al. Prove2Me: An Open Collaborative Platform for Scaling Math Formalization. arXiv:2608.28433, 2026.
\item Anthropic. Formalizing Fermat's Last Theorem. anthropic.com/research, 2026.
\item N. Carlini. Building a C Compiler with a Team of Parallel Claudes. Anthropic Engineering, 2026.
\item L. Lamport. Specifying Systems: The TLA+ Language and Tools for Hardware and Software Engineers. Addison-Wesley, 2002.
\item P. Hoffman, J. Schlyter. The DNS-Based Authentication of Named Entities (DANE) Transport Layer Security (TLS) Protocol: TLSA. RFC 6698, 2012.
\item P. Wouters. DNS-Based Authentication of Named Entities (DANE) Bindings for OpenPGP. RFC 7929, 2016.
\item M. Kleppmann. Designing Data-Intensive Applications. O'Reilly, 2017. (Chapters on replication, consensus and leases; recommended background reading.)
\end{enumerate}
\end{english}

\end{document}
